% Options for packages loaded elsewhere
\PassOptionsToPackage{unicode}{hyperref}
\PassOptionsToPackage{hyphens}{url}
\PassOptionsToPackage{dvipsnames,svgnames,x11names}{xcolor}
%
\documentclass[
  10pt,
]{article}

\usepackage{amsmath,amssymb}
\usepackage{iftex}
\ifPDFTeX
  \usepackage[T1]{fontenc}
  \usepackage[utf8]{inputenc}
  \usepackage{textcomp} % provide euro and other symbols
\else % if luatex or xetex
  \usepackage{unicode-math}
  \defaultfontfeatures{Scale=MatchLowercase}
  \defaultfontfeatures[\rmfamily]{Ligatures=TeX,Scale=1}
\fi
\usepackage{lmodern}
\ifPDFTeX\else  
    % xetex/luatex font selection
    \setmainfont[]{Arial}
\fi
% Use upquote if available, for straight quotes in verbatim environments
\IfFileExists{upquote.sty}{\usepackage{upquote}}{}
\IfFileExists{microtype.sty}{% use microtype if available
  \usepackage[]{microtype}
  \UseMicrotypeSet[protrusion]{basicmath} % disable protrusion for tt fonts
}{}
\makeatletter
\@ifundefined{KOMAClassName}{% if non-KOMA class
  \IfFileExists{parskip.sty}{%
    \usepackage{parskip}
  }{% else
    \setlength{\parindent}{0pt}
    \setlength{\parskip}{6pt plus 2pt minus 1pt}}
}{% if KOMA class
  \KOMAoptions{parskip=half}}
\makeatother
\usepackage{xcolor}
\usepackage[a4paper]{geometry}
\setlength{\emergencystretch}{3em} % prevent overfull lines
\setcounter{secnumdepth}{5}
% Make \paragraph and \subparagraph free-standing
\makeatletter
\ifx\paragraph\undefined\else
  \let\oldparagraph\paragraph
  \renewcommand{\paragraph}{
    \@ifstar
      \xxxParagraphStar
      \xxxParagraphNoStar
  }
  \newcommand{\xxxParagraphStar}[1]{\oldparagraph*{#1}\mbox{}}
  \newcommand{\xxxParagraphNoStar}[1]{\oldparagraph{#1}\mbox{}}
\fi
\ifx\subparagraph\undefined\else
  \let\oldsubparagraph\subparagraph
  \renewcommand{\subparagraph}{
    \@ifstar
      \xxxSubParagraphStar
      \xxxSubParagraphNoStar
  }
  \newcommand{\xxxSubParagraphStar}[1]{\oldsubparagraph*{#1}\mbox{}}
  \newcommand{\xxxSubParagraphNoStar}[1]{\oldsubparagraph{#1}\mbox{}}
\fi
\makeatother


\providecommand{\tightlist}{%
  \setlength{\itemsep}{0pt}\setlength{\parskip}{0pt}}\usepackage{longtable,booktabs,array}
\usepackage{calc} % for calculating minipage widths
% Correct order of tables after \paragraph or \subparagraph
\usepackage{etoolbox}
\makeatletter
\patchcmd\longtable{\par}{\if@noskipsec\mbox{}\fi\par}{}{}
\makeatother
% Allow footnotes in longtable head/foot
\IfFileExists{footnotehyper.sty}{\usepackage{footnotehyper}}{\usepackage{footnote}}
\makesavenoteenv{longtable}
\usepackage{graphicx}
\makeatletter
\newsavebox\pandoc@box
\newcommand*\pandocbounded[1]{% scales image to fit in text height/width
  \sbox\pandoc@box{#1}%
  \Gscale@div\@tempa{\textheight}{\dimexpr\ht\pandoc@box+\dp\pandoc@box\relax}%
  \Gscale@div\@tempb{\linewidth}{\wd\pandoc@box}%
  \ifdim\@tempb\p@<\@tempa\p@\let\@tempa\@tempb\fi% select the smaller of both
  \ifdim\@tempa\p@<\p@\scalebox{\@tempa}{\usebox\pandoc@box}%
  \else\usebox{\pandoc@box}%
  \fi%
}
% Set default figure placement to htbp
\def\fps@figure{htbp}
\makeatother

\makeatletter
\@ifpackageloaded{caption}{}{\usepackage{caption}}
\AtBeginDocument{%
\ifdefined\contentsname
  \renewcommand*\contentsname{Table of contents}
\else
  \newcommand\contentsname{Table of contents}
\fi
\ifdefined\listfigurename
  \renewcommand*\listfigurename{List of Figures}
\else
  \newcommand\listfigurename{List of Figures}
\fi
\ifdefined\listtablename
  \renewcommand*\listtablename{List of Tables}
\else
  \newcommand\listtablename{List of Tables}
\fi
\ifdefined\figurename
  \renewcommand*\figurename{Figure}
\else
  \newcommand\figurename{Figure}
\fi
\ifdefined\tablename
  \renewcommand*\tablename{Table}
\else
  \newcommand\tablename{Table}
\fi
}
\@ifpackageloaded{float}{}{\usepackage{float}}
\floatstyle{ruled}
\@ifundefined{c@chapter}{\newfloat{codelisting}{h}{lop}}{\newfloat{codelisting}{h}{lop}[chapter]}
\floatname{codelisting}{Listing}
\newcommand*\listoflistings{\listof{codelisting}{List of Listings}}
\makeatother
\makeatletter
\makeatother
\makeatletter
\@ifpackageloaded{caption}{}{\usepackage{caption}}
\@ifpackageloaded{subcaption}{}{\usepackage{subcaption}}
\makeatother

\usepackage{bookmark}

\IfFileExists{xurl.sty}{\usepackage{xurl}}{} % add URL line breaks if available
\urlstyle{same} % disable monospaced font for URLs
\hypersetup{
  pdftitle={04\_Mutants},
  pdfauthor={Henry Janse van Rensburg},
  colorlinks=true,
  linkcolor={blue},
  filecolor={Maroon},
  citecolor={Blue},
  urlcolor={Blue},
  pdfcreator={LaTeX via pandoc}}


\title{04\_Mutants}
\author{Henry Janse van Rensburg}
\date{2026-11-15}

\begin{document}
\maketitle

\renewcommand*\contentsname{Table of contents}
{
\hypersetup{linkcolor=}
\setcounter{tocdepth}{2}
\tableofcontents
}

\section{Background}\label{background}

This document analyses mutant feedback experiments testing plant growth
responses to BXplus- and BXminus-conditioned soil microbiomes.

For the *mmf1* mutant screen, three experiments are shown individually
in detail:

- E-III: *mmf1* repetition 1

- E-IV: *mmf1* repetition 2

- E-V: *mmf1* repetition 3

A final combined *mmf1* analysis is then performed using pre-combined
tabs in `Dataset4`, which include five experiments in total.

For the ETI mutant screen, one experiment is shown individually in
detail:

- E-VI: ETI mutants

A final combined ETI analysis is then performed using pre-combined tabs
in `Dataset4`, which include four experiments in total.

The *mmf1* complementation lines are analysed separately as:

- E-X: *mmf1* complementation lines

\section{Experiment overview}\label{experiment-overview}

\begin{longtable}[]{@{}
  >{\raggedright\arraybackslash}p{(\linewidth - 6\tabcolsep) * \real{0.2500}}
  >{\raggedright\arraybackslash}p{(\linewidth - 6\tabcolsep) * \real{0.2500}}
  >{\raggedright\arraybackslash}p{(\linewidth - 6\tabcolsep) * \real{0.2500}}
  >{\raggedright\arraybackslash}p{(\linewidth - 6\tabcolsep) * \real{0.2500}}@{}}
\toprule\noalign{}
\begin{minipage}[b]{\linewidth}\raggedright
Experiment
\end{minipage} & \begin{minipage}[b]{\linewidth}\raggedright
Description
\end{minipage} & \begin{minipage}[b]{\linewidth}\raggedright
Soil batch
\end{minipage} & \begin{minipage}[b]{\linewidth}\raggedright
Notes
\end{minipage} \\
\midrule\noalign{}
\endhead
\bottomrule\noalign{}
\endlastfoot
E-III & \emph{mmf1} repetition 1 & BS04 & Growth \\
E-IV & \emph{mmf1} repetition 2 & BS04 & Growth, microbiota,
transcriptome \\
E-V & \emph{mmf1} repetition 3 & BS09 & Growth \\
E-VI & ETI mutants & BS07 & Growth \\
E-VII & ETI mutants & BS07 & Growth \\
E-VIII & ETI mutants & BS11 & Growth, transcriptome \\
E-IX & ETI mutants & BS10 & Growth \\
E-X & \emph{mmf1} complementation lines & BS10 & Growth \\
\end{longtable}

\section{Setup}\label{setup}

\begin{verbatim}
 [1] "ExpIII_mmf1_weight"         "ExpIII_mmf1_area"          
 [3] "ExpIV_mmf1_weight"          "ExpIV_mmf1_area"           
 [5] "ExpV_mmf1_weight"           "ExpV_mmf1_area"            
 [7] "ExpVI_ETI_weight"           "ExpVI_ETI_area"            
 [9] "ExpVII_complimentation"     "ExpVII_ETI_combined_weight"
[11] "ExpVII_ETI_combined_area"   "mmf1_combined_weight"      
[13] "mmf1_combined_area"        
\end{verbatim}

\subsection{Helper function}\label{helper-function}

\newpage

\section{\texorpdfstring{\textbf{Data
import}}{Data import}}\label{data-import}

\textbf{Individual mmf1 experiments shown in detail}

\subsection{Factoring and cleaning
data}\label{factoring-and-cleaning-data}

\newpage

\section{\texorpdfstring{\textbf{mmf1 mutant
screen}}{mmf1 mutant screen}}\label{mmf1-mutant-screen}

\subsection{\texorpdfstring{\textbf{Experiment III: mmf1-1 on
BS04}}{Experiment III: mmf1-1 on BS04}}\label{experiment-iii-mmf1-1-on-bs04}

\subsubsection{\texorpdfstring{\textbf{Replicates}}{Replicates}}\label{replicates}

\begin{longtable}[]{@{}lrr@{}}
\caption{Number of replicates for fresh weight (E-III)}\tabularnewline
\toprule\noalign{}
& WT & mmf1-1 \\
\midrule\noalign{}
\endfirsthead
\toprule\noalign{}
& WT & mmf1-1 \\
\midrule\noalign{}
\endhead
\bottomrule\noalign{}
\endlastfoot
minus & 14 & 14 \\
plus & 15 & 15 \\
\end{longtable}

\begin{longtable}[]{@{}lrr@{}}
\caption{Number of replicates for rosette area (E-III)}\tabularnewline
\toprule\noalign{}
& WT & mmf1-1 \\
\midrule\noalign{}
\endfirsthead
\toprule\noalign{}
& WT & mmf1-1 \\
\midrule\noalign{}
\endhead
\bottomrule\noalign{}
\endlastfoot
minus & 14 & 14 \\
plus & 15 & 15 \\
\end{longtable}

\subsubsection{\texorpdfstring{\textbf{Fresh
weight}}{Fresh weight}}\label{fresh-weight}

\begin{longtable}[]{@{}lrrr@{}}
\caption{Mean fresh weight and Δ(plus-minus) by genotype
(E-III)}\tabularnewline
\toprule\noalign{}
line & minus & plus & delta\_plus\_minus \\
\midrule\noalign{}
\endfirsthead
\toprule\noalign{}
line & minus & plus & delta\_plus\_minus \\
\midrule\noalign{}
\endhead
\bottomrule\noalign{}
\endlastfoot
WT & 0.0431436 & 0.0591273 & 0.0159838 \\
mmf1-1 & 0.0446614 & 0.0523180 & 0.0076566 \\
\end{longtable}

\begin{figure}

\centering{

\pandocbounded{\includegraphics[keepaspectratio]{04_Mutant_screening_files/figure-pdf/fig-mmf1-fw-EIII-1.pdf}}

}

\caption{\label{fig-mmf1-fw-EIII}Fresh weight of WT and mmf1-1 in
E-III.}

\end{figure}%

\subsubsection{Rosette area}\label{rosette-area}

\begin{longtable}[]{@{}lrrr@{}}
\caption{Mean rosette area and Δ(plus-minus) by genotype
(E-III)}\tabularnewline
\toprule\noalign{}
line & minus & plus & delta\_plus\_minus \\
\midrule\noalign{}
\endfirsthead
\toprule\noalign{}
line & minus & plus & delta\_plus\_minus \\
\midrule\noalign{}
\endhead
\bottomrule\noalign{}
\endlastfoot
WT & 432457.3 & 614835.9 & 182378.65 \\
mmf1-1 & 461693.9 & 530039.2 & 68345.34 \\
\end{longtable}

\begin{figure}

\centering{

\pandocbounded{\includegraphics[keepaspectratio]{04_Mutant_screening_files/figure-pdf/fig-mmf1-area-EIII-1.pdf}}

}

\caption{\label{fig-mmf1-area-EIII}Rosette area of WT and mmf1-1 in
E-III.}

\end{figure}%

\subsubsection{Statistics}\label{statistics}

\begin{longtable}[]{@{}lrrrrr@{}}
\caption{ANOVA table for fresh weight (E-III)}\tabularnewline
\toprule\noalign{}
& Df & Sum Sq & Mean Sq & F value & Pr(\textgreater F) \\
\midrule\noalign{}
\endfirsthead
\toprule\noalign{}
& Df & Sum Sq & Mean Sq & F value & Pr(\textgreater F) \\
\midrule\noalign{}
\endhead
\bottomrule\noalign{}
\endlastfoot
line & 1 & 0.0001128 & 0.0001128 & 0.5779622 & 0.4504188 \\
BX & 1 & 0.0020235 & 0.0020235 & 10.3665983 & 0.0021736 \\
line:BX & 1 & 0.0002511 & 0.0002511 & 1.2862520 & 0.2617495 \\
Residuals & 54 & 0.0105404 & 0.0001952 & NA & NA \\
\end{longtable}

\begin{longtable}[]{@{}lrrrrr@{}}
\caption{ANOVA table for rosette area (E-III)}\tabularnewline
\toprule\noalign{}
& Df & Sum Sq & Mean Sq & F value & Pr(\textgreater F) \\
\midrule\noalign{}
\endfirsthead
\toprule\noalign{}
& Df & Sum Sq & Mean Sq & F value & Pr(\textgreater F) \\
\midrule\noalign{}
\endhead
\bottomrule\noalign{}
\endlastfoot
line & 1 & 1.283010e+10 & 12830104212 & 0.6702395 & 0.4165665 \\
BX & 1 & 2.276057e+11 & 227605673691 & 11.8900284 & 0.0011014 \\
line:BX & 1 & 4.708198e+10 & 47081980430 & 2.4595436 & 0.1226545 \\
Residuals & 54 & 1.033699e+12 & 19142567756 & NA & NA \\
\end{longtable}

\begin{longtable}[]{@{}lrrl@{}}
\caption{Per-line BX effect with FDR correction for fresh weight
(E-III)}\tabularnewline
\toprule\noalign{}
line & p\_value & p\_adj & significance \\
\midrule\noalign{}
\endfirsthead
\toprule\noalign{}
line & p\_value & p\_adj & significance \\
\midrule\noalign{}
\endhead
\bottomrule\noalign{}
\endlastfoot
WT & 0.0159009 & 0.0318019 & * \\
mmf1-1 & 0.0608325 & 0.0608325 & ns \\
\end{longtable}

\begin{longtable}[]{@{}lrrl@{}}
\caption{Per-line BX effect with FDR correction for rosette area
(E-III)}\tabularnewline
\toprule\noalign{}
line & p\_value & p\_adj & significance \\
\midrule\noalign{}
\endfirsthead
\toprule\noalign{}
line & p\_value & p\_adj & significance \\
\midrule\noalign{}
\endhead
\bottomrule\noalign{}
\endlastfoot
WT & 0.0039416 & 0.0078831 & ** \\
mmf1-1 & 0.1323721 & 0.1323721 & ns \\
\end{longtable}

\newpage

\subsection{\texorpdfstring{\textbf{Experiment IV: mmf1-1 and mmf1-2 on
BS04}}{Experiment IV: mmf1-1 and mmf1-2 on BS04}}\label{experiment-iv-mmf1-1-and-mmf1-2-on-bs04}

This experiment was also used for microbiota and transcriptome analyses.

\subsubsection{\texorpdfstring{\textbf{Replicates}}{Replicates}}\label{replicates-1}

\begin{longtable}[]{@{}lrrr@{}}
\caption{Number of replicates for fresh weight (E-IV)}\tabularnewline
\toprule\noalign{}
& WT & mmf1-1 & mmf1-2 \\
\midrule\noalign{}
\endfirsthead
\toprule\noalign{}
& WT & mmf1-1 & mmf1-2 \\
\midrule\noalign{}
\endhead
\bottomrule\noalign{}
\endlastfoot
minus & 15 & 15 & 15 \\
plus & 15 & 14 & 14 \\
\end{longtable}

\begin{longtable}[]{@{}lrrr@{}}
\caption{Number of replicates for rosette area (E-IV)}\tabularnewline
\toprule\noalign{}
& WT & mmf1-1 & mmf1-2 \\
\midrule\noalign{}
\endfirsthead
\toprule\noalign{}
& WT & mmf1-1 & mmf1-2 \\
\midrule\noalign{}
\endhead
\bottomrule\noalign{}
\endlastfoot
minus & 15 & 15 & 15 \\
plus & 15 & 14 & 14 \\
\end{longtable}

\subsubsection{Fresh weight}\label{fresh-weight-1}

\begin{longtable}[]{@{}lrrr@{}}
\caption{Mean fresh weight and Δ(plus-minus) by genotype
(E-IV)}\tabularnewline
\toprule\noalign{}
line & minus & plus & delta\_plus\_minus \\
\midrule\noalign{}
\endfirsthead
\toprule\noalign{}
line & minus & plus & delta\_plus\_minus \\
\midrule\noalign{}
\endhead
\bottomrule\noalign{}
\endlastfoot
WT & 0.2168520 & 0.2402087 & 0.0233567 \\
mmf1-1 & 0.1553347 & 0.1597164 & 0.0043818 \\
mmf1-2 & 0.1991680 & 0.2085900 & 0.0094220 \\
\end{longtable}

\begin{figure}

\centering{

\pandocbounded{\includegraphics[keepaspectratio]{04_Mutant_screening_files/figure-pdf/fig-mmf1-fw-EIV-1.pdf}}

}

\caption{\label{fig-mmf1-fw-EIV}Figure 4b (fresh weight) \textbar{} MMF1
is involved in growth feedback responses to the BXplus microbiome. (a)
Representative rosettes of 6-week-old Col-0 (WT) and two independent
mmf1 T-DNA insertion mutants (AT3G44630) grown on BXminus and BXplus
soil microbiomes (experiment IV, Table S1). Scale bar, 2 cm. (b) Fresh
weight (left) and rosette area (right; in pixels) after 6 weeks of
growth for the genotypes in (a). Boxplots show the median (centre line)
and the 25th and 75th percentiles (box bounds); whiskers extend to the
most extreme value within 1.5 × the interquartile range of the nearest
box bound; points represent individual plants grown one per pot,
i.e.~biological replicates (sample sizes (n) are given below each
boxplot). Text above each genotype indicates the observed difference in
means between conditions, expressed as Δmean = BXplus − BXminus.
Statistical significance was assessed within each genotype by one-way
ANOVA of BXplus versus BXminus (with two groups, equivalent to a
two-sided t-test with pooled variance); P values were
Benjamini--Hochberg-adjusted across the three genotypes and the adjusted
P values are shown. The findings were reproduced in five additional
experiments (E-III, E-V, E-VIII, E-IX and E-X) on four separately
conditioned soil batches (see Extended Data Fig. 5). Dataset 4.}

\end{figure}%

\subsubsection{Rosette area}\label{rosette-area-1}

\begin{longtable}[]{@{}lrrr@{}}
\caption{Mean rosette area and Δ(plus-minus) by genotype
(E-IV)}\tabularnewline
\toprule\noalign{}
line & minus & plus & delta\_plus\_minus \\
\midrule\noalign{}
\endfirsthead
\toprule\noalign{}
line & minus & plus & delta\_plus\_minus \\
\midrule\noalign{}
\endhead
\bottomrule\noalign{}
\endlastfoot
WT & 428236.5 & 513806.5 & 85570.067 \\
mmf1-1 & 322590.7 & 348724.5 & 26133.833 \\
mmf1-2 & 446723.8 & 442869.1 & -3854.729 \\
\end{longtable}

\begin{figure}

\centering{

\pandocbounded{\includegraphics[keepaspectratio]{04_Mutant_screening_files/figure-pdf/fig-mmf1-area-EIV-1.pdf}}

}

\caption{\label{fig-mmf1-area-EIV}Figure 4b (rosette area) \textbar{}
MMF1 is involved in growth feedback responses to the BXplus microbiome.
(a) Representative rosettes of 6-week-old Col-0 (WT) and two independent
mmf1 T-DNA insertion mutants (AT3G44630) grown on BXminus and BXplus
soil microbiomes (experiment IV, Table S1). Scale bar, 2 cm. (b) Fresh
weight (left) and rosette area (right; in pixels) after 6 weeks of
growth for the genotypes in (a). Boxplots show the median (centre line)
and the 25th and 75th percentiles (box bounds); whiskers extend to the
most extreme value within 1.5 × the interquartile range of the nearest
box bound; points represent individual plants grown one per pot,
i.e.~biological replicates (sample sizes (n) are given below each
boxplot). Text above each genotype indicates the observed difference in
means between conditions, expressed as Δmean = BXplus − BXminus.
Statistical significance was assessed within each genotype by one-way
ANOVA of BXplus versus BXminus (with two groups, equivalent to a
two-sided t-test with pooled variance); P values were
Benjamini--Hochberg-adjusted across the three genotypes and the adjusted
P values are shown. The findings were reproduced in five additional
experiments (E-III, E-V, E-VIII, E-IX and E-X) on four separately
conditioned soil batches (see Extended Data Fig. 5). Dataset 4.}

\end{figure}%

\newpage

\subsubsection{Statistics}\label{statistics-1}

\begin{longtable}[]{@{}lrrrrr@{}}
\caption{ANOVA table for fresh weight (E-IV)}\tabularnewline
\toprule\noalign{}
& Df & Sum Sq & Mean Sq & F value & Pr(\textgreater F) \\
\midrule\noalign{}
\endfirsthead
\toprule\noalign{}
& Df & Sum Sq & Mean Sq & F value & Pr(\textgreater F) \\
\midrule\noalign{}
\endhead
\bottomrule\noalign{}
\endlastfoot
line & 2 & 0.0764945 & 0.0382473 & 89.646602 & 0.0000000 \\
BX & 1 & 0.0034435 & 0.0034435 & 8.071188 & 0.0056715 \\
line:BX & 2 & 0.0014299 & 0.0007149 & 1.675694 & 0.1935235 \\
Residuals & 82 & 0.0349849 & 0.0004266 & NA & NA \\
\end{longtable}

\begin{longtable}[]{@{}lrrrrr@{}}
\caption{ANOVA table for rosette area (E-IV)}\tabularnewline
\toprule\noalign{}
& Df & Sum Sq & Mean Sq & F value & Pr(\textgreater F) \\
\midrule\noalign{}
\endfirsthead
\toprule\noalign{}
& Df & Sum Sq & Mean Sq & F value & Pr(\textgreater F) \\
\midrule\noalign{}
\endhead
\bottomrule\noalign{}
\endlastfoot
line & 2 & 304039816844 & 152019908422 & 21.457669 & 0.0000000 \\
BX & 1 & 29340301711 & 29340301711 & 4.141395 & 0.0450772 \\
line:BX & 2 & 30629767064 & 15314883532 & 2.161702 & 0.1216458 \\
Residuals & 82 & 580940661226 & 7084642210 & NA & NA \\
\end{longtable}

\begin{longtable}[]{@{}lrrl@{}}
\caption{Per-line BX effect with FDR correction for fresh weight
(E-IV)}\tabularnewline
\toprule\noalign{}
line & p\_value & p\_adj & significance \\
\midrule\noalign{}
\endfirsthead
\toprule\noalign{}
line & p\_value & p\_adj & significance \\
\midrule\noalign{}
\endhead
\bottomrule\noalign{}
\endlastfoot
WT & 0.0146263 & 0.0438788 & * \\
mmf1-1 & 0.4793053 & 0.4793053 & ns \\
mmf1-2 & 0.2152555 & 0.3228833 & ns \\
\end{longtable}

\begin{longtable}[]{@{}lrrl@{}}
\caption{Per-line BX effect with FDR correction for rosette area
(E-IV)}\tabularnewline
\toprule\noalign{}
line & p\_value & p\_adj & significance \\
\midrule\noalign{}
\endfirsthead
\toprule\noalign{}
line & p\_value & p\_adj & significance \\
\midrule\noalign{}
\endhead
\bottomrule\noalign{}
\endlastfoot
WT & 0.0112422 & 0.0337265 & * \\
mmf1-1 & 0.2127660 & 0.3191489 & ns \\
mmf1-2 & 0.9210433 & 0.9210433 & ns \\
\end{longtable}

\begin{longtable}[]{@{}
  >{\raggedright\arraybackslash}p{(\linewidth - 10\tabcolsep) * \real{0.4103}}
  >{\raggedleft\arraybackslash}p{(\linewidth - 10\tabcolsep) * \real{0.1410}}
  >{\raggedleft\arraybackslash}p{(\linewidth - 10\tabcolsep) * \real{0.1282}}
  >{\raggedleft\arraybackslash}p{(\linewidth - 10\tabcolsep) * \real{0.0385}}
  >{\raggedleft\arraybackslash}p{(\linewidth - 10\tabcolsep) * \real{0.1538}}
  >{\raggedleft\arraybackslash}p{(\linewidth - 10\tabcolsep) * \real{0.1282}}@{}}
\caption{Pairwise Sidak-adjusted comparisons for fresh weight in
E-IV}\tabularnewline
\toprule\noalign{}
\begin{minipage}[b]{\linewidth}\raggedright
contrast
\end{minipage} & \begin{minipage}[b]{\linewidth}\raggedleft
estimate
\end{minipage} & \begin{minipage}[b]{\linewidth}\raggedleft
SE
\end{minipage} & \begin{minipage}[b]{\linewidth}\raggedleft
df
\end{minipage} & \begin{minipage}[b]{\linewidth}\raggedleft
t.ratio
\end{minipage} & \begin{minipage}[b]{\linewidth}\raggedleft
p.value
\end{minipage} \\
\midrule\noalign{}
\endfirsthead
\toprule\noalign{}
\begin{minipage}[b]{\linewidth}\raggedright
contrast
\end{minipage} & \begin{minipage}[b]{\linewidth}\raggedleft
estimate
\end{minipage} & \begin{minipage}[b]{\linewidth}\raggedleft
SE
\end{minipage} & \begin{minipage}[b]{\linewidth}\raggedleft
df
\end{minipage} & \begin{minipage}[b]{\linewidth}\raggedleft
t.ratio
\end{minipage} & \begin{minipage}[b]{\linewidth}\raggedleft
p.value
\end{minipage} \\
\midrule\noalign{}
\endhead
\bottomrule\noalign{}
\endlastfoot
WT minus - (mmf1-1 minus) & 0.0615173 & 0.0075423 & 82 & 8.1563310 &
0.0000000 \\
WT minus - (mmf1-2 minus) & 0.0176840 & 0.0075423 & 82 & 2.3446491 &
0.2778011 \\
WT minus - WT plus & -0.0233567 & 0.0075423 & 82 & -3.0967647 &
0.0394266 \\
WT minus - (mmf1-1 plus) & 0.0571356 & 0.0076758 & 82 & 7.4436156 &
0.0000000 \\
WT minus - (mmf1-2 plus) & 0.0082620 & 0.0076758 & 82 & 1.0763724 &
0.9934635 \\
(mmf1-1 minus) - (mmf1-2 minus) & -0.0438333 & 0.0075423 & 82 &
-5.8116820 & 0.0000017 \\
(mmf1-1 minus) - WT plus & -0.0848740 & 0.0075423 & 82 & -11.2530957 &
0.0000000 \\
(mmf1-1 minus) - (mmf1-1 plus) & -0.0043818 & 0.0076758 & 82 &
-0.5708554 & 0.9999968 \\
(mmf1-1 minus) - (mmf1-2 plus) & -0.0532553 & 0.0076758 & 82 &
-6.9380986 & 0.0000000 \\
(mmf1-2 minus) - WT plus & -0.0410407 & 0.0075423 & 82 & -5.4414137 &
0.0000080 \\
(mmf1-2 minus) - (mmf1-1 plus) & 0.0394516 & 0.0076758 & 82 & 5.1397461
& 0.0000275 \\
(mmf1-2 minus) - (mmf1-2 plus) & -0.0094220 & 0.0076758 & 82 &
-1.2274971 & 0.9773497 \\
WT plus - (mmf1-1 plus) & 0.0804922 & 0.0076758 & 82 & 10.4865194 &
0.0000000 \\
WT plus - (mmf1-2 plus) & 0.0316187 & 0.0076758 & 82 & 4.1192762 &
0.0013529 \\
(mmf1-1 plus) - (mmf1-2 plus) & -0.0488736 & 0.0078070 & 82 & -6.2602231
& 0.0000003 \\
\end{longtable}

\begin{longtable}[]{@{}
  >{\raggedright\arraybackslash}p{(\linewidth - 16\tabcolsep) * \real{0.0455}}
  >{\raggedright\arraybackslash}p{(\linewidth - 16\tabcolsep) * \real{0.1061}}
  >{\raggedright\arraybackslash}p{(\linewidth - 16\tabcolsep) * \real{0.0909}}
  >{\raggedleft\arraybackslash}p{(\linewidth - 16\tabcolsep) * \real{0.1515}}
  >{\raggedleft\arraybackslash}p{(\linewidth - 16\tabcolsep) * \real{0.1515}}
  >{\raggedleft\arraybackslash}p{(\linewidth - 16\tabcolsep) * \real{0.0455}}
  >{\raggedleft\arraybackslash}p{(\linewidth - 16\tabcolsep) * \real{0.1515}}
  >{\raggedleft\arraybackslash}p{(\linewidth - 16\tabcolsep) * \real{0.1515}}
  >{\raggedright\arraybackslash}p{(\linewidth - 16\tabcolsep) * \real{0.1061}}@{}}
\caption{Multiple comparison groupings (Sidak-adjusted) for fresh weight
in E-IV}\tabularnewline
\toprule\noalign{}
\begin{minipage}[b]{\linewidth}\raggedright
\end{minipage} & \begin{minipage}[b]{\linewidth}\raggedright
line
\end{minipage} & \begin{minipage}[b]{\linewidth}\raggedright
BX
\end{minipage} & \begin{minipage}[b]{\linewidth}\raggedleft
emmean
\end{minipage} & \begin{minipage}[b]{\linewidth}\raggedleft
SE
\end{minipage} & \begin{minipage}[b]{\linewidth}\raggedleft
df
\end{minipage} & \begin{minipage}[b]{\linewidth}\raggedleft
lower.CL
\end{minipage} & \begin{minipage}[b]{\linewidth}\raggedleft
upper.CL
\end{minipage} & \begin{minipage}[b]{\linewidth}\raggedright
.group
\end{minipage} \\
\midrule\noalign{}
\endfirsthead
\toprule\noalign{}
\begin{minipage}[b]{\linewidth}\raggedright
\end{minipage} & \begin{minipage}[b]{\linewidth}\raggedright
line
\end{minipage} & \begin{minipage}[b]{\linewidth}\raggedright
BX
\end{minipage} & \begin{minipage}[b]{\linewidth}\raggedleft
emmean
\end{minipage} & \begin{minipage}[b]{\linewidth}\raggedleft
SE
\end{minipage} & \begin{minipage}[b]{\linewidth}\raggedleft
df
\end{minipage} & \begin{minipage}[b]{\linewidth}\raggedleft
lower.CL
\end{minipage} & \begin{minipage}[b]{\linewidth}\raggedleft
upper.CL
\end{minipage} & \begin{minipage}[b]{\linewidth}\raggedright
.group
\end{minipage} \\
\midrule\noalign{}
\endhead
\bottomrule\noalign{}
\endlastfoot
2 & mmf1-1 & minus & 0.1553347 & 0.0053332 & 82 & 0.1409561 & 0.1697132
& a \\
5 & mmf1-1 & plus & 0.1597164 & 0.0055204 & 82 & 0.1448332 & 0.1745997 &
a \\
3 & mmf1-2 & minus & 0.1991680 & 0.0053332 & 82 & 0.1847894 & 0.2135466
& b \\
6 & mmf1-2 & plus & 0.2085900 & 0.0055204 & 82 & 0.1937068 & 0.2234732 &
b \\
1 & WT & minus & 0.2168520 & 0.0053332 & 82 & 0.2024734 & 0.2312306 &
b \\
4 & WT & plus & 0.2402087 & 0.0053332 & 82 & 0.2258301 & 0.2545872 &
c \\
\end{longtable}

\begin{longtable}[]{@{}
  >{\raggedright\arraybackslash}p{(\linewidth - 10\tabcolsep) * \real{0.4156}}
  >{\raggedleft\arraybackslash}p{(\linewidth - 10\tabcolsep) * \real{0.1558}}
  >{\raggedleft\arraybackslash}p{(\linewidth - 10\tabcolsep) * \real{0.1169}}
  >{\raggedleft\arraybackslash}p{(\linewidth - 10\tabcolsep) * \real{0.0390}}
  >{\raggedleft\arraybackslash}p{(\linewidth - 10\tabcolsep) * \real{0.1429}}
  >{\raggedleft\arraybackslash}p{(\linewidth - 10\tabcolsep) * \real{0.1299}}@{}}
\caption{Pairwise Sidak-adjusted comparisons for rosette area in
E-IV}\tabularnewline
\toprule\noalign{}
\begin{minipage}[b]{\linewidth}\raggedright
contrast
\end{minipage} & \begin{minipage}[b]{\linewidth}\raggedleft
estimate
\end{minipage} & \begin{minipage}[b]{\linewidth}\raggedleft
SE
\end{minipage} & \begin{minipage}[b]{\linewidth}\raggedleft
df
\end{minipage} & \begin{minipage}[b]{\linewidth}\raggedleft
t.ratio
\end{minipage} & \begin{minipage}[b]{\linewidth}\raggedleft
p.value
\end{minipage} \\
\midrule\noalign{}
\endfirsthead
\toprule\noalign{}
\begin{minipage}[b]{\linewidth}\raggedright
contrast
\end{minipage} & \begin{minipage}[b]{\linewidth}\raggedleft
estimate
\end{minipage} & \begin{minipage}[b]{\linewidth}\raggedleft
SE
\end{minipage} & \begin{minipage}[b]{\linewidth}\raggedleft
df
\end{minipage} & \begin{minipage}[b]{\linewidth}\raggedleft
t.ratio
\end{minipage} & \begin{minipage}[b]{\linewidth}\raggedleft
p.value
\end{minipage} \\
\midrule\noalign{}
\endhead
\bottomrule\noalign{}
\endlastfoot
WT minus - (mmf1-1 minus) & 105645.800 & 30734.65 & 82 & 3.4373512 &
0.0137880 \\
WT minus - (mmf1-2 minus) & -18487.333 & 30734.65 & 82 & -0.6015143 &
0.9999935 \\
WT minus - WT plus & -85570.067 & 30734.65 & 82 & -2.7841558 &
0.0953958 \\
WT minus - (mmf1-1 plus) & 79511.967 & 31278.67 & 82 & 2.5420506 &
0.1769952 \\
WT minus - (mmf1-2 plus) & -14632.605 & 31278.67 & 82 & -0.4678141 &
0.9999998 \\
(mmf1-1 minus) - (mmf1-2 minus) & -124133.133 & 30734.65 & 82 &
-4.0388655 & 0.0018027 \\
(mmf1-1 minus) - WT plus & -191215.867 & 30734.65 & 82 & -6.2215070 &
0.0000003 \\
(mmf1-1 minus) - (mmf1-1 plus) & -26133.833 & 31278.67 & 82 & -0.8355161
& 0.9995941 \\
(mmf1-1 minus) - (mmf1-2 plus) & -120278.405 & 31278.67 & 82 &
-3.8453807 & 0.0035469 \\
(mmf1-2 minus) - WT plus & -67082.733 & 30734.65 & 82 & -2.1826416 &
0.3853149 \\
(mmf1-2 minus) - (mmf1-1 plus) & 97999.300 & 31278.67 & 82 & 3.1331029 &
0.0353885 \\
(mmf1-2 minus) - (mmf1-2 plus) & 3854.729 & 31278.67 & 82 & 0.1232382 &
1.0000000 \\
WT plus - (mmf1-1 plus) & 165082.033 & 31278.67 & 82 & 5.2777826 &
0.0000157 \\
WT plus - (mmf1-2 plus) & 70937.462 & 31278.67 & 82 & 2.2679179 &
0.3260543 \\
(mmf1-1 plus) - (mmf1-2 plus) & -94144.571 & 31813.39 & 82 & -2.9592751
& 0.0587703 \\
\end{longtable}

\begin{longtable}[]{@{}lllrrrrrl@{}}
\caption{Multiple comparison groupings (Sidak-adjusted) for rosette area
in E-IV}\tabularnewline
\toprule\noalign{}
& line & BX & emmean & SE & df & lower.CL & upper.CL & .group \\
\midrule\noalign{}
\endfirsthead
\toprule\noalign{}
& line & BX & emmean & SE & df & lower.CL & upper.CL & .group \\
\midrule\noalign{}
\endhead
\bottomrule\noalign{}
\endlastfoot
2 & mmf1-1 & minus & 322590.7 & 21732.68 & 82 & 263998.2 & 381183.1 &
a \\
5 & mmf1-1 & plus & 348724.5 & 22495.46 & 82 & 288075.5 & 409373.5 &
ab \\
1 & WT & minus & 428236.5 & 21732.68 & 82 & 369644.0 & 486828.9 & bc \\
6 & mmf1-2 & plus & 442869.1 & 22495.46 & 82 & 382220.1 & 503518.0 &
bc \\
3 & mmf1-2 & minus & 446723.8 & 21732.68 & 82 & 388131.3 & 505316.3 &
c \\
4 & WT & plus & 513806.5 & 21732.68 & 82 & 455214.1 & 572399.0 & c \\
\end{longtable}

\newpage

\subsection{\texorpdfstring{\textbf{Experiment V: mmf1-1 and mmf1-2 on
BS09}}{Experiment V: mmf1-1 and mmf1-2 on BS09}}\label{experiment-v-mmf1-1-and-mmf1-2-on-bs09}

\subsubsection{\texorpdfstring{\textbf{Replicates}}{Replicates}}\label{replicates-2}

\begin{longtable}[]{@{}lrrr@{}}
\caption{Number of replicates for fresh weight (E-V)}\tabularnewline
\toprule\noalign{}
& WT & mmf1-1 & mmf1-2 \\
\midrule\noalign{}
\endfirsthead
\toprule\noalign{}
& WT & mmf1-1 & mmf1-2 \\
\midrule\noalign{}
\endhead
\bottomrule\noalign{}
\endlastfoot
minus & 14 & 10 & 13 \\
plus & 13 & 10 & 14 \\
\end{longtable}

\begin{longtable}[]{@{}lrrr@{}}
\caption{Number of replicates for rosette area (E-V)}\tabularnewline
\toprule\noalign{}
& WT & mmf1-1 & mmf1-2 \\
\midrule\noalign{}
\endfirsthead
\toprule\noalign{}
& WT & mmf1-1 & mmf1-2 \\
\midrule\noalign{}
\endhead
\bottomrule\noalign{}
\endlastfoot
minus & 14 & 10 & 13 \\
plus & 13 & 10 & 14 \\
\end{longtable}

\subsubsection{Fresh weight}\label{fresh-weight-2}

\begin{longtable}[]{@{}lrrr@{}}
\caption{Mean fresh weight and Δ(plus-minus) by genotype
(E-V)}\tabularnewline
\toprule\noalign{}
line & minus & plus & delta\_plus\_minus \\
\midrule\noalign{}
\endfirsthead
\toprule\noalign{}
line & minus & plus & delta\_plus\_minus \\
\midrule\noalign{}
\endhead
\bottomrule\noalign{}
\endlastfoot
WT & 0.2346614 & 0.2879615 & 0.0533001 \\
mmf1-1 & 0.0616500 & 0.0725600 & 0.0109100 \\
mmf1-2 & 0.1240231 & 0.1252071 & 0.0011841 \\
\end{longtable}

\begin{figure}

\centering{

\pandocbounded{\includegraphics[keepaspectratio]{04_Mutant_screening_files/figure-pdf/fig-mmf1-fw-EV-1.pdf}}

}

\caption{\label{fig-mmf1-fw-EV}Fresh weight of WT, mmf1-1 and mmf1-2 in
E-V.}

\end{figure}%

\subsubsection{Rosette area}\label{rosette-area-2}

\begin{longtable}[]{@{}lrrr@{}}
\caption{Mean rosette area and Δ(plus-minus) by genotype
(E-V)}\tabularnewline
\toprule\noalign{}
line & minus & plus & delta\_plus\_minus \\
\midrule\noalign{}
\endfirsthead
\toprule\noalign{}
line & minus & plus & delta\_plus\_minus \\
\midrule\noalign{}
\endhead
\bottomrule\noalign{}
\endlastfoot
WT & 1503286.6 & 1760205.7 & 256919.12 \\
mmf1-1 & 477854.8 & 556068.8 & 78214.00 \\
mmf1-2 & 858764.6 & 857408.3 & -1356.33 \\
\end{longtable}

\begin{figure}

\centering{

\pandocbounded{\includegraphics[keepaspectratio]{04_Mutant_screening_files/figure-pdf/fig-mmf1-area-EV-1.pdf}}

}

\caption{\label{fig-mmf1-area-EV}Rosette area of WT, mmf1-1 and mmf1-2
in E-V.}

\end{figure}%

\subsubsection{Statistics}\label{statistics-2}

\begin{longtable}[]{@{}lrrrrr@{}}
\caption{ANOVA table for fresh weight (E-V)}\tabularnewline
\toprule\noalign{}
& Df & Sum Sq & Mean Sq & F value & Pr(\textgreater F) \\
\midrule\noalign{}
\endfirsthead
\toprule\noalign{}
& Df & Sum Sq & Mean Sq & F value & Pr(\textgreater F) \\
\midrule\noalign{}
\endhead
\bottomrule\noalign{}
\endlastfoot
line & 2 & 0.4779630 & 0.2389815 & 204.209912 & 0.0000000 \\
BX & 1 & 0.0096273 & 0.0096273 & 8.226538 & 0.0054934 \\
line:BX & 2 & 0.0101271 & 0.0050635 & 4.326796 & 0.0170293 \\
Residuals & 68 & 0.0795786 & 0.0011703 & NA & NA \\
\end{longtable}

\begin{longtable}[]{@{}lrrrrr@{}}
\caption{ANOVA table for rosette area (E-V)}\tabularnewline
\toprule\noalign{}
& Df & Sum Sq & Mean Sq & F value & Pr(\textgreater F) \\
\midrule\noalign{}
\endfirsthead
\toprule\noalign{}
& Df & Sum Sq & Mean Sq & F value & Pr(\textgreater F) \\
\midrule\noalign{}
\endhead
\bottomrule\noalign{}
\endlastfoot
line & 2 & 1.566506e+13 & 7.832531e+12 & 195.131941 & 0.0000000 \\
BX & 1 & 2.417528e+11 & 2.417528e+11 & 6.022790 & 0.0166883 \\
line:BX & 2 & 2.337858e+11 & 1.168929e+11 & 2.912154 & 0.0611680 \\
Residuals & 68 & 2.729497e+12 & 4.013967e+10 & NA & NA \\
\end{longtable}

\begin{longtable}[]{@{}lrrl@{}}
\caption{Per-line BX effect with FDR correction for fresh weight
(E-V)}\tabularnewline
\toprule\noalign{}
line & p\_value & p\_adj & significance \\
\midrule\noalign{}
\endfirsthead
\toprule\noalign{}
line & p\_value & p\_adj & significance \\
\midrule\noalign{}
\endhead
\bottomrule\noalign{}
\endlastfoot
WT & 0.0027711 & 0.0083132 & ** \\
mmf1-1 & 0.1689862 & 0.2534794 & ns \\
mmf1-2 & 0.9310226 & 0.9310226 & ns \\
\end{longtable}

\begin{longtable}[]{@{}lrrl@{}}
\caption{Per-line BX effect with FDR correction for rosette area
(E-V)}\tabularnewline
\toprule\noalign{}
line & p\_value & p\_adj & significance \\
\midrule\noalign{}
\endfirsthead
\toprule\noalign{}
line & p\_value & p\_adj & significance \\
\midrule\noalign{}
\endhead
\bottomrule\noalign{}
\endlastfoot
WT & 0.0088133 & 0.0264399 & * \\
mmf1-1 & 0.1452810 & 0.2179215 & ns \\
mmf1-2 & 0.9868226 & 0.9868226 & ns \\
\end{longtable}

\newpage

\subsection{\texorpdfstring{\textbf{mmf1 complementation lines
(E-X)}}{mmf1 complementation lines (E-X)}}\label{mmf1-complementation-lines-e-x}

\subsubsection{\texorpdfstring{\textbf{Replicates}}{Replicates}}\label{replicates-3}

\begin{longtable}[]{@{}lrrrr@{}}
\caption{Number of replicates for complementation lines
(E-X)}\tabularnewline
\toprule\noalign{}
& WT & mmf1-2 & mmf1-2*Nos4-4 & mmf1-2*Nos5-1 \\
\midrule\noalign{}
\endfirsthead
\toprule\noalign{}
& WT & mmf1-2 & mmf1-2*Nos4-4 & mmf1-2*Nos5-1 \\
\midrule\noalign{}
\endhead
\bottomrule\noalign{}
\endlastfoot
minus & 13 & 15 & 14 & 15 \\
plus & 14 & 13 & 14 & 14 \\
\end{longtable}

\subsubsection{Statistics}\label{statistics-3}

\begin{longtable}[]{@{}lrrrrr@{}}
\caption{ANOVA table for fresh weight (E-X)}\tabularnewline
\toprule\noalign{}
& Df & Sum Sq & Mean Sq & F value & Pr(\textgreater F) \\
\midrule\noalign{}
\endfirsthead
\toprule\noalign{}
& Df & Sum Sq & Mean Sq & F value & Pr(\textgreater F) \\
\midrule\noalign{}
\endhead
\bottomrule\noalign{}
\endlastfoot
line & 3 & 0.5209019 & 0.1736340 & 56.751765 & 0.0000000 \\
BX & 1 & 0.0160766 & 0.0160766 & 5.254581 & 0.0239029 \\
line:BX & 3 & 0.0634982 & 0.0211661 & 6.918069 & 0.0002718 \\
Residuals & 104 & 0.3181916 & 0.0030595 & NA & NA \\
\end{longtable}

\begin{longtable}[]{@{}lrrrrr@{}}
\caption{ANOVA table for rosette area (E-X)}\tabularnewline
\toprule\noalign{}
& Df & Sum Sq & Mean Sq & F value & Pr(\textgreater F) \\
\midrule\noalign{}
\endfirsthead
\toprule\noalign{}
& Df & Sum Sq & Mean Sq & F value & Pr(\textgreater F) \\
\midrule\noalign{}
\endhead
\bottomrule\noalign{}
\endlastfoot
line & 3 & 1.687669e+12 & 562556426788 & 47.091592 & 0.0000000 \\
BX & 1 & 5.684074e+10 & 56840737126 & 4.758138 & 0.0314126 \\
line:BX & 3 & 2.089085e+11 & 69636177463 & 5.829244 & 0.0010123 \\
Residuals & 104 & 1.242385e+12 & 11946005611 & NA & NA \\
\end{longtable}

\begin{longtable}[]{@{}lrrrrr@{}}
\caption{ANOVA table for dry weight (E-X)}\tabularnewline
\toprule\noalign{}
& Df & Sum Sq & Mean Sq & F value & Pr(\textgreater F) \\
\midrule\noalign{}
\endfirsthead
\toprule\noalign{}
& Df & Sum Sq & Mean Sq & F value & Pr(\textgreater F) \\
\midrule\noalign{}
\endhead
\bottomrule\noalign{}
\endlastfoot
line & 3 & 11708.7096 & 3902.90319 & 57.407099 & 0.000000 \\
BX & 1 & 409.7172 & 409.71725 & 6.026457 & 0.015752 \\
line:BX & 3 & 1644.7840 & 548.26134 & 8.064277 & 0.000070 \\
Residuals & 104 & 7070.5877 & 67.98642 & NA & NA \\
\end{longtable}

\begin{longtable}[]{@{}lrrl@{}}
\caption{Per-line BX effect with FDR correction for fresh weight
(E-X)}\tabularnewline
\toprule\noalign{}
line & p\_value & p\_adj & significance \\
\midrule\noalign{}
\endfirsthead
\toprule\noalign{}
line & p\_value & p\_adj & significance \\
\midrule\noalign{}
\endhead
\bottomrule\noalign{}
\endlastfoot
WT & 0.0102704 & 0.0205408 & * \\
mmf1-2 & 0.0343600 & 0.0458133 & * \\
mmf1-2*Nos4-4 & 0.0028343 & 0.0113371 & * \\
mmf1-2*Nos5-1 & 0.1715184 & 0.1715184 & ns \\
\end{longtable}

\begin{longtable}[]{@{}lrrl@{}}
\caption{Per-line BX effect with FDR correction for rosette area
(E-X)}\tabularnewline
\toprule\noalign{}
line & p\_value & p\_adj & significance \\
\midrule\noalign{}
\endfirsthead
\toprule\noalign{}
line & p\_value & p\_adj & significance \\
\midrule\noalign{}
\endhead
\bottomrule\noalign{}
\endlastfoot
WT & 0.0025974 & 0.0103895 & * \\
mmf1-2 & 0.1207216 & 0.1609621 & ns \\
mmf1-2*Nos4-4 & 0.0093994 & 0.0187987 & * \\
mmf1-2*Nos5-1 & 0.6697610 & 0.6697610 & ns \\
\end{longtable}

\begin{longtable}[]{@{}lrrl@{}}
\caption{Per-line BX effect with FDR correction for dry weight
(E-X)}\tabularnewline
\toprule\noalign{}
line & p\_value & p\_adj & significance \\
\midrule\noalign{}
\endfirsthead
\toprule\noalign{}
line & p\_value & p\_adj & significance \\
\midrule\noalign{}
\endhead
\bottomrule\noalign{}
\endlastfoot
WT & 0.0074153 & 0.0148306 & * \\
mmf1-2 & 0.0220332 & 0.0293776 & * \\
mmf1-2*Nos4-4 & 0.0011682 & 0.0046727 & ** \\
mmf1-2*Nos5-1 & 0.1600319 & 0.1600319 & ns \\
\end{longtable}

\subsubsection{Fresh weight}\label{fresh-weight-3}

\begin{figure}

\centering{

\pandocbounded{\includegraphics[keepaspectratio]{04_Mutant_screening_files/figure-pdf/fig-compl-fw-EX-1.pdf}}

}

\caption{\label{fig-compl-fw-EX}Fresh weight of mmf1 complementation
lines in E-X.}

\end{figure}%

\subsubsection{Dry weight}\label{dry-weight}

\begin{figure}

\centering{

\pandocbounded{\includegraphics[keepaspectratio]{04_Mutant_screening_files/figure-pdf/fig-compl-dw-EX-1.pdf}}

}

\caption{\label{fig-compl-dw-EX}Dry weight of mmf1 complementation lines
in E-X.}

\end{figure}%

\subsubsection{Rosette area}\label{rosette-area-3}

\begin{figure}

\centering{

\pandocbounded{\includegraphics[keepaspectratio]{04_Mutant_screening_files/figure-pdf/fig-compl-area-EX-1.pdf}}

}

\caption{\label{fig-compl-area-EX}Rosette area of mmf1 complementation
lines in E-X.}

\end{figure}%

\newpage

\section{\texorpdfstring{\textbf{ETI mutant
screen}}{ETI mutant screen}}\label{eti-mutant-screen}

\subsection{\texorpdfstring{\textbf{Experiment VI: ETI mutants on
BS07}}{Experiment VI: ETI mutants on BS07}}\label{experiment-vi-eti-mutants-on-bs07}

\subsubsection{\texorpdfstring{\textbf{Replicates}}{Replicates}}\label{replicates-4}

\begin{longtable}[]{@{}lrrrr@{}}
\caption{Number of replicates - weight (E-VI)}\tabularnewline
\toprule\noalign{}
& WT & eds1 & pad4 & sag101 \\
\midrule\noalign{}
\endfirsthead
\toprule\noalign{}
& WT & eds1 & pad4 & sag101 \\
\midrule\noalign{}
\endhead
\bottomrule\noalign{}
\endlastfoot
minus & 14 & 14 & 11 & 11 \\
plus & 14 & 13 & 14 & 12 \\
\end{longtable}

\begin{longtable}[]{@{}lrrrr@{}}
\caption{Number of replicates - plant area (E-VI)}\tabularnewline
\toprule\noalign{}
& WT & eds1 & pad4 & sag101 \\
\midrule\noalign{}
\endfirsthead
\toprule\noalign{}
& WT & eds1 & pad4 & sag101 \\
\midrule\noalign{}
\endhead
\bottomrule\noalign{}
\endlastfoot
minus & 14 & 14 & 11 & 11 \\
plus & 14 & 13 & 14 & 12 \\
\end{longtable}

\subsubsection{Statistics}\label{statistics-4}

\begin{longtable}[]{@{}lrrrrr@{}}
\caption{ANOVA table for fresh weight (E-VI)}\tabularnewline
\toprule\noalign{}
& Df & Sum Sq & Mean Sq & F value & Pr(\textgreater F) \\
\midrule\noalign{}
\endfirsthead
\toprule\noalign{}
& Df & Sum Sq & Mean Sq & F value & Pr(\textgreater F) \\
\midrule\noalign{}
\endhead
\bottomrule\noalign{}
\endlastfoot
line & 3 & 6442.0308 & 2147.3436 & 0.4960015 & 0.6859221 \\
BX & 1 & 101029.3326 & 101029.3326 & 23.3361373 & 0.0000052 \\
line:BX & 3 & 342.7104 & 114.2368 & 0.0263868 & 0.9941739 \\
Residuals & 95 & 411284.2869 & 4329.3083 & NA & NA \\
\end{longtable}

\begin{longtable}[]{@{}lrrrrr@{}}
\caption{ANOVA table for dry weight (E-VI)}\tabularnewline
\toprule\noalign{}
& Df & Sum Sq & Mean Sq & F value & Pr(\textgreater F) \\
\midrule\noalign{}
\endfirsthead
\toprule\noalign{}
& Df & Sum Sq & Mean Sq & F value & Pr(\textgreater F) \\
\midrule\noalign{}
\endhead
\bottomrule\noalign{}
\endlastfoot
line & 3 & 29.266438 & 9.755479 & 0.2857264 & 0.8355970 \\
BX & 1 & 891.529645 & 891.529645 & 26.1118475 & 0.0000017 \\
line:BX & 3 & 4.622734 & 1.540911 & 0.0451315 & 0.9871912 \\
Residuals & 95 & 3243.558929 & 34.142726 & NA & NA \\
\end{longtable}

\begin{longtable}[]{@{}lrrrrr@{}}
\caption{ANOVA table for rosette area (E-VI)}\tabularnewline
\toprule\noalign{}
& Df & Sum Sq & Mean Sq & F value & Pr(\textgreater F) \\
\midrule\noalign{}
\endfirsthead
\toprule\noalign{}
& Df & Sum Sq & Mean Sq & F value & Pr(\textgreater F) \\
\midrule\noalign{}
\endhead
\bottomrule\noalign{}
\endlastfoot
line & 3 & 1.706198e+12 & 5.687326e+11 & 1.9316359 & 0.1296917 \\
BX & 1 & 5.775147e+12 & 5.775147e+12 & 19.6146312 & 0.0000253 \\
line:BX & 3 & 7.012305e+10 & 2.337435e+10 & 0.0793883 & 0.9710410 \\
Residuals & 95 & 2.797090e+13 & 2.944306e+11 & NA & NA \\
\end{longtable}

\begin{longtable}[]{@{}lrrl@{}}
\caption{Per-line BX effect with FDR correction for fresh weight
(E-VI)}\tabularnewline
\toprule\noalign{}
line & p\_value & p\_adj & significance \\
\midrule\noalign{}
\endfirsthead
\toprule\noalign{}
line & p\_value & p\_adj & significance \\
\midrule\noalign{}
\endhead
\bottomrule\noalign{}
\endlastfoot
WT & 0.0029714 & 0.0118855 & * \\
eds1 & 0.0310466 & 0.0456622 & * \\
pad4 & 0.0525278 & 0.0525278 & ns \\
sag101 & 0.0342466 & 0.0456622 & * \\
\end{longtable}

\begin{longtable}[]{@{}lrrl@{}}
\caption{Per-line BX effect with FDR correction for dry weight
(E-VI)}\tabularnewline
\toprule\noalign{}
line & p\_value & p\_adj & significance \\
\midrule\noalign{}
\endfirsthead
\toprule\noalign{}
line & p\_value & p\_adj & significance \\
\midrule\noalign{}
\endhead
\bottomrule\noalign{}
\endlastfoot
WT & 0.0018612 & 0.0074447 & ** \\
eds1 & 0.0261247 & 0.0353938 & * \\
pad4 & 0.0373417 & 0.0373417 & * \\
sag101 & 0.0265453 & 0.0353938 & * \\
\end{longtable}

\begin{longtable}[]{@{}lrrl@{}}
\caption{Per-line BX effect with FDR correction for rosette area
(E-VI)}\tabularnewline
\toprule\noalign{}
line & p\_value & p\_adj & significance \\
\midrule\noalign{}
\endfirsthead
\toprule\noalign{}
line & p\_value & p\_adj & significance \\
\midrule\noalign{}
\endhead
\bottomrule\noalign{}
\endlastfoot
WT & 0.0086979 & 0.0347917 & * \\
eds1 & 0.0256198 & 0.0512396 & ns \\
pad4 & 0.0803324 & 0.0803324 & ns \\
sag101 & 0.0555717 & 0.0740956 & ns \\
\end{longtable}

\subsubsection{Dry weight}\label{dry-weight-1}

\begin{figure}

\centering{

\pandocbounded{\includegraphics[keepaspectratio]{04_Mutant_screening_files/figure-pdf/fig-eti-dw-EVI-1.pdf}}

}

\caption{\label{fig-eti-dw-EVI}Dry weight of ETI mutants in E-VI.}

\end{figure}%

\subsubsection{Fresh weight}\label{fresh-weight-4}

\begin{figure}

\centering{

\pandocbounded{\includegraphics[keepaspectratio]{04_Mutant_screening_files/figure-pdf/fig-eti-fw-EVI-1.pdf}}

}

\caption{\label{fig-eti-fw-EVI}Fresh weight of ETI mutants in E-VI.}

\end{figure}%

\subsubsection{Rosette area}\label{rosette-area-4}

\begin{figure}

\centering{

\pandocbounded{\includegraphics[keepaspectratio]{04_Mutant_screening_files/figure-pdf/fig-eti-area-EVI-1.pdf}}

}

\caption{\label{fig-eti-area-EVI}Rosette area of ETI mutants in E-VI.}

\end{figure}%

\newpage

\subsection{\texorpdfstring{\textbf{Combined ETI analysis across E-VI to
E-IX}}{Combined ETI analysis across E-VI to E-IX}}\label{combined-eti-analysis-across-e-vi-to-e-ix}

\subsubsection{\texorpdfstring{\textbf{Replicates}}{Replicates}}\label{replicates-5}

\begin{longtable}[]{@{}lrrrr@{}}
\caption{Number of replicates - weight (ETI combined, E-VI to
E-IX)}\tabularnewline
\toprule\noalign{}
& WT & eds1 & pad4 & sag101 \\
\midrule\noalign{}
\endfirsthead
\toprule\noalign{}
& WT & eds1 & pad4 & sag101 \\
\midrule\noalign{}
\endhead
\bottomrule\noalign{}
\endlastfoot
minus & 52 & 52 & 44 & 48 \\
plus & 53 & 52 & 48 & 53 \\
\end{longtable}

\begin{longtable}[]{@{}lrrrr@{}}
\caption{Number of replicates - plant area (ETI combined, E-VI to
E-IX)}\tabularnewline
\toprule\noalign{}
& WT & eds1 & pad4 & sag101 \\
\midrule\noalign{}
\endfirsthead
\toprule\noalign{}
& WT & eds1 & pad4 & sag101 \\
\midrule\noalign{}
\endhead
\bottomrule\noalign{}
\endlastfoot
minus & 52 & 52 & 44 & 48 \\
plus & 53 & 52 & 48 & 53 \\
\end{longtable}

\subsubsection{LMM and EMM: fresh
weight}\label{lmm-and-emm-fresh-weight}

\begin{verbatim}
Linear mixed model fit by REML. t-tests use Satterthwaite's method [
lmerModLmerTest]
Formula: formula_lmm
   Data: data

REML criterion at convergence: 4787.3

Scaled residuals: 
    Min      1Q  Median      3Q     Max 
-3.2434 -0.5362 -0.0469  0.4914  3.9864 

Random effects:
 Groups     Name        Variance Std.Dev.
 experiment (Intercept) 11867    108.93  
 Residual                9865     99.32  
Number of obs: 402, groups:  experiment, 4

Fixed effects:
                  Estimate Std. Error      df t value Pr(>|t|)    
(Intercept)        197.696     56.183   3.341   3.519 0.032801 *  
BXplus              69.738     19.394 391.006   3.596 0.000365 ***
lineeds1            12.848     19.486 391.005   0.659 0.510044    
linepad4           -28.713     20.373 391.018  -1.409 0.159526    
linesag101         -12.417     19.886 391.004  -0.624 0.532709    
BXplus:lineeds1    -40.495     27.491 391.005  -1.473 0.141552    
BXplus:linepad4    -12.436     28.396 391.007  -0.438 0.661677    
BXplus:linesag101  -13.303     27.721 391.009  -0.480 0.631559    
---
Signif. codes:  0 '***' 0.001 '**' 0.01 '*' 0.05 '.' 0.1 ' ' 1

Correlation of Fixed Effects:
            (Intr) BXplus linds1 linpd4 lns101 BXpl:1 BXpl:4
BXplus      -0.174                                          
lineeds1    -0.173  0.502                                   
linepad4    -0.166  0.480  0.479                            
linesag101  -0.170  0.492  0.490  0.469                     
BXpls:lnds1  0.123 -0.706 -0.708 -0.339 -0.347              
BXpls:lnpd4  0.119 -0.683 -0.343 -0.716 -0.336  0.482       
BXpls:ln101  0.122 -0.700 -0.351 -0.336 -0.717  0.494  0.478
\end{verbatim}

\begin{longtable}[]{@{}
  >{\raggedright\arraybackslash}p{(\linewidth - 16\tabcolsep) * \real{0.1444}}
  >{\raggedright\arraybackslash}p{(\linewidth - 16\tabcolsep) * \real{0.0778}}
  >{\raggedleft\arraybackslash}p{(\linewidth - 16\tabcolsep) * \real{0.2000}}
  >{\raggedleft\arraybackslash}p{(\linewidth - 16\tabcolsep) * \real{0.0778}}
  >{\raggedleft\arraybackslash}p{(\linewidth - 16\tabcolsep) * \real{0.1000}}
  >{\raggedleft\arraybackslash}p{(\linewidth - 16\tabcolsep) * \real{0.1000}}
  >{\raggedleft\arraybackslash}p{(\linewidth - 16\tabcolsep) * \real{0.1000}}
  >{\raggedleft\arraybackslash}p{(\linewidth - 16\tabcolsep) * \real{0.1000}}
  >{\raggedleft\arraybackslash}p{(\linewidth - 16\tabcolsep) * \real{0.1000}}@{}}
\caption{Within-line ΔEMM = BXplus - BXminus for ETI fresh
weight}\tabularnewline
\toprule\noalign{}
\begin{minipage}[b]{\linewidth}\raggedright
contrast
\end{minipage} & \begin{minipage}[b]{\linewidth}\raggedright
line
\end{minipage} & \begin{minipage}[b]{\linewidth}\raggedleft
effect\_plus\_minus
\end{minipage} & \begin{minipage}[b]{\linewidth}\raggedleft
SE
\end{minipage} & \begin{minipage}[b]{\linewidth}\raggedleft
df
\end{minipage} & \begin{minipage}[b]{\linewidth}\raggedleft
lower.CL
\end{minipage} & \begin{minipage}[b]{\linewidth}\raggedleft
upper.CL
\end{minipage} & \begin{minipage}[b]{\linewidth}\raggedleft
t
\end{minipage} & \begin{minipage}[b]{\linewidth}\raggedleft
p
\end{minipage} \\
\midrule\noalign{}
\endfirsthead
\toprule\noalign{}
\begin{minipage}[b]{\linewidth}\raggedright
contrast
\end{minipage} & \begin{minipage}[b]{\linewidth}\raggedright
line
\end{minipage} & \begin{minipage}[b]{\linewidth}\raggedleft
effect\_plus\_minus
\end{minipage} & \begin{minipage}[b]{\linewidth}\raggedleft
SE
\end{minipage} & \begin{minipage}[b]{\linewidth}\raggedleft
df
\end{minipage} & \begin{minipage}[b]{\linewidth}\raggedleft
lower.CL
\end{minipage} & \begin{minipage}[b]{\linewidth}\raggedleft
upper.CL
\end{minipage} & \begin{minipage}[b]{\linewidth}\raggedleft
t
\end{minipage} & \begin{minipage}[b]{\linewidth}\raggedleft
p
\end{minipage} \\
\midrule\noalign{}
\endhead
\bottomrule\noalign{}
\endlastfoot
plus - minus & WT & 69.738 & 19.394 & 391.0056 & 31.607 & 107.868 &
3.595754 & 0.000365 \\
plus - minus & eds1 & 29.243 & 19.482 & 391.0025 & -9.060 & 67.546 &
1.500990 & 0.134000 \\
plus - minus & pad4 & 57.302 & 20.736 & 391.0038 & 16.534 & 98.070 &
2.763436 & 0.005990 \\
plus - minus & sag101 & 56.434 & 19.795 & 391.0031 & 17.517 & 95.351 &
2.851014 & 0.004590 \\
\end{longtable}

Extended Data Figure 6c (fresh weight) \textbar{} Growth feedback in a
BXplus microbiome remains intact in effector-triggered immunity
signalling mutants. (a) Phenotypes of 6-week-old WT, and mutants
(eds1-12, pad4-1 and sag101-3) affected in effector-triggered immunity
signalling pathways grown on the BXminus and BXplus soils. Scale bar, 2
cm. (b) Z-score--transformed fresh weight and rosette area of WT,
eds1-12, pad4-1, and sag101-3 plants across four independent experiments
illustrate the positive growth feedback on the BXplus soil microbiome.
Values were z-score--transformed within experiment to account for
between-experiment variation. Each point is an individual plant grown in
its own pot, i.e.~a biological replicate; number of replicates is given
below each plot. Boxplots show the median (centre line) and the 25th and
75th percentiles (box bounds); whiskers extend to the most extreme value
within 1.5 × the interquartile range of the nearest box bound, and all
plants are plotted as points. The statistics above the plots are
estimated marginal mean (EMM) contrasts of the BX effect within each
genotype (ΔEMM = BXplus − BXminus), computed from a linear mixed-effects
model with experiment as a random intercept to account for variation
between experiments. Each ΔEMM is a two-sided t-test on the estimated
marginal means (emmeans). ΔEMM is therefore expressed in the units of
the measurement, not in z-score units. (c) Summary of the same linear
mixed-effects model in (b). For each genotype, the BX effect is reported
as the EMM contrast Δ(BXplus − BXminus) with unadjusted two-sided P
values. The table also reports interaction contrasts relative to WT (ΔΔ
vs WT = Δ\_mutant − Δ\_WT), which test whether the BX response of each
mutant differs from that of WT; these three comparisons form one family
per measure and are two-sided t-tests with Holm adjusted P values. Data
are from Dataset 4.

\begin{longtable}[]{@{}
  >{\raggedright\arraybackslash}p{(\linewidth - 16\tabcolsep) * \real{0.1596}}
  >{\raggedright\arraybackslash}p{(\linewidth - 16\tabcolsep) * \real{0.1809}}
  >{\raggedleft\arraybackslash}p{(\linewidth - 16\tabcolsep) * \real{0.1170}}
  >{\raggedleft\arraybackslash}p{(\linewidth - 16\tabcolsep) * \real{0.0745}}
  >{\raggedleft\arraybackslash}p{(\linewidth - 16\tabcolsep) * \real{0.0957}}
  >{\raggedleft\arraybackslash}p{(\linewidth - 16\tabcolsep) * \real{0.0957}}
  >{\raggedleft\arraybackslash}p{(\linewidth - 16\tabcolsep) * \real{0.0957}}
  >{\raggedleft\arraybackslash}p{(\linewidth - 16\tabcolsep) * \real{0.1170}}
  >{\raggedleft\arraybackslash}p{(\linewidth - 16\tabcolsep) * \real{0.0638}}@{}}
\caption{Difference-in-differences relative to WT for ETI fresh
weight}\tabularnewline
\toprule\noalign{}
\begin{minipage}[b]{\linewidth}\raggedright
BX\_revpairwise
\end{minipage} & \begin{minipage}[b]{\linewidth}\raggedright
line\_trt.vs.ctrl
\end{minipage} & \begin{minipage}[b]{\linewidth}\raggedleft
diff\_vs\_WT
\end{minipage} & \begin{minipage}[b]{\linewidth}\raggedleft
SE
\end{minipage} & \begin{minipage}[b]{\linewidth}\raggedleft
df
\end{minipage} & \begin{minipage}[b]{\linewidth}\raggedleft
lower.CL
\end{minipage} & \begin{minipage}[b]{\linewidth}\raggedleft
upper.CL
\end{minipage} & \begin{minipage}[b]{\linewidth}\raggedleft
t
\end{minipage} & \begin{minipage}[b]{\linewidth}\raggedleft
p\_adj
\end{minipage} \\
\midrule\noalign{}
\endfirsthead
\toprule\noalign{}
\begin{minipage}[b]{\linewidth}\raggedright
BX\_revpairwise
\end{minipage} & \begin{minipage}[b]{\linewidth}\raggedright
line\_trt.vs.ctrl
\end{minipage} & \begin{minipage}[b]{\linewidth}\raggedleft
diff\_vs\_WT
\end{minipage} & \begin{minipage}[b]{\linewidth}\raggedleft
SE
\end{minipage} & \begin{minipage}[b]{\linewidth}\raggedleft
df
\end{minipage} & \begin{minipage}[b]{\linewidth}\raggedleft
lower.CL
\end{minipage} & \begin{minipage}[b]{\linewidth}\raggedleft
upper.CL
\end{minipage} & \begin{minipage}[b]{\linewidth}\raggedleft
t
\end{minipage} & \begin{minipage}[b]{\linewidth}\raggedleft
p\_adj
\end{minipage} \\
\midrule\noalign{}
\endhead
\bottomrule\noalign{}
\endlastfoot
plus - minus & eds1 - WT & -40.495 & 27.491 & 391.0047 & -106.594 &
25.603 & -1.4730085 & 0.425 \\
plus - minus & pad4 - WT & -12.436 & 28.397 & 391.0065 & -80.710 &
55.839 & -0.4379305 & 1.000 \\
plus - minus & sag101 - WT & -13.303 & 27.721 & 391.0084 & -79.953 &
53.346 & -0.4799077 & 1.000 \\
\end{longtable}

Extended Data Figure 6c (fresh weight) \textbar{} Growth feedback in a
BXplus microbiome remains intact in effector-triggered immunity
signalling mutants. (a) Phenotypes of 6-week-old WT, and mutants
(eds1-12, pad4-1 and sag101-3) affected in effector-triggered immunity
signalling pathways grown on the BXminus and BXplus soils. Scale bar, 2
cm. (b) Z-score--transformed fresh weight and rosette area of WT,
eds1-12, pad4-1, and sag101-3 plants across four independent experiments
illustrate the positive growth feedback on the BXplus soil microbiome.
Values were z-score--transformed within experiment to account for
between-experiment variation. Each point is an individual plant grown in
its own pot, i.e.~a biological replicate; number of replicates is given
below each plot. Boxplots show the median (centre line) and the 25th and
75th percentiles (box bounds); whiskers extend to the most extreme value
within 1.5 × the interquartile range of the nearest box bound, and all
plants are plotted as points. The statistics above the plots are
estimated marginal mean (EMM) contrasts of the BX effect within each
genotype (ΔEMM = BXplus − BXminus), computed from a linear mixed-effects
model with experiment as a random intercept to account for variation
between experiments. Each ΔEMM is a two-sided t-test on the estimated
marginal means (emmeans). ΔEMM is therefore expressed in the units of
the measurement, not in z-score units. (c) Summary of the same linear
mixed-effects model in (b). For each genotype, the BX effect is reported
as the EMM contrast Δ(BXplus − BXminus) with unadjusted two-sided P
values. The table also reports interaction contrasts relative to WT (ΔΔ
vs WT = Δ\_mutant − Δ\_WT), which test whether the BX response of each
mutant differs from that of WT; these three comparisons form one family
per measure and are two-sided t-tests with Holm adjusted P values. Data
are from Dataset 4.

\begin{longtable}[]{@{}
  >{\raggedright\arraybackslash}p{(\linewidth - 16\tabcolsep) * \real{0.1500}}
  >{\raggedright\arraybackslash}p{(\linewidth - 16\tabcolsep) * \real{0.0750}}
  >{\raggedleft\arraybackslash}p{(\linewidth - 16\tabcolsep) * \real{0.1125}}
  >{\raggedleft\arraybackslash}p{(\linewidth - 16\tabcolsep) * \real{0.0875}}
  >{\raggedleft\arraybackslash}p{(\linewidth - 16\tabcolsep) * \real{0.1125}}
  >{\raggedleft\arraybackslash}p{(\linewidth - 16\tabcolsep) * \real{0.1125}}
  >{\raggedleft\arraybackslash}p{(\linewidth - 16\tabcolsep) * \real{0.1125}}
  >{\raggedleft\arraybackslash}p{(\linewidth - 16\tabcolsep) * \real{0.1375}}
  >{\raggedleft\arraybackslash}p{(\linewidth - 16\tabcolsep) * \real{0.1000}}@{}}
\caption{Mutant vs WT contrasts within each BX condition for ETI fresh
weight (Holm-adjusted)}\tabularnewline
\toprule\noalign{}
\begin{minipage}[b]{\linewidth}\raggedright
contrast
\end{minipage} & \begin{minipage}[b]{\linewidth}\raggedright
BX
\end{minipage} & \begin{minipage}[b]{\linewidth}\raggedleft
estimate
\end{minipage} & \begin{minipage}[b]{\linewidth}\raggedleft
SE
\end{minipage} & \begin{minipage}[b]{\linewidth}\raggedleft
df
\end{minipage} & \begin{minipage}[b]{\linewidth}\raggedleft
lower.CL
\end{minipage} & \begin{minipage}[b]{\linewidth}\raggedleft
upper.CL
\end{minipage} & \begin{minipage}[b]{\linewidth}\raggedleft
t.ratio
\end{minipage} & \begin{minipage}[b]{\linewidth}\raggedleft
p.value
\end{minipage} \\
\midrule\noalign{}
\endfirsthead
\toprule\noalign{}
\begin{minipage}[b]{\linewidth}\raggedright
contrast
\end{minipage} & \begin{minipage}[b]{\linewidth}\raggedright
BX
\end{minipage} & \begin{minipage}[b]{\linewidth}\raggedleft
estimate
\end{minipage} & \begin{minipage}[b]{\linewidth}\raggedleft
SE
\end{minipage} & \begin{minipage}[b]{\linewidth}\raggedleft
df
\end{minipage} & \begin{minipage}[b]{\linewidth}\raggedleft
lower.CL
\end{minipage} & \begin{minipage}[b]{\linewidth}\raggedleft
upper.CL
\end{minipage} & \begin{minipage}[b]{\linewidth}\raggedleft
t.ratio
\end{minipage} & \begin{minipage}[b]{\linewidth}\raggedleft
p.value
\end{minipage} \\
\midrule\noalign{}
\endhead
\bottomrule\noalign{}
\endlastfoot
eds1 - WT & minus & 12.848 & 19.486 & 391.0048 & -34.002 & 59.699 &
0.6593706 & 1.000 \\
pad4 - WT & minus & -28.713 & 20.374 & 391.0179 & -77.698 & 20.272 &
-1.4093294 & 0.479 \\
sag101 - WT & minus & -12.417 & 19.886 & 391.0037 & -60.229 & 35.395 &
-0.6244287 & 1.000 \\
eds1 - WT & plus & -27.647 & 19.404 & 391.0121 & -74.300 & 19.007 &
-1.4247924 & 0.310 \\
pad4 - WT & plus & -41.149 & 19.834 & 391.0291 & -88.837 & 6.540 &
-2.0746140 & 0.116 \\
sag101 - WT & plus & -25.721 & 19.330 & 391.0248 & -72.196 & 20.754 &
-1.3306221 & 0.310 \\
\end{longtable}

Extended Data Figure 6c (fresh weight) \textbar{} Growth feedback in a
BXplus microbiome remains intact in effector-triggered immunity
signalling mutants. (a) Phenotypes of 6-week-old WT, and mutants
(eds1-12, pad4-1 and sag101-3) affected in effector-triggered immunity
signalling pathways grown on the BXminus and BXplus soils. Scale bar, 2
cm. (b) Z-score--transformed fresh weight and rosette area of WT,
eds1-12, pad4-1, and sag101-3 plants across four independent experiments
illustrate the positive growth feedback on the BXplus soil microbiome.
Values were z-score--transformed within experiment to account for
between-experiment variation. Each point is an individual plant grown in
its own pot, i.e.~a biological replicate; number of replicates is given
below each plot. Boxplots show the median (centre line) and the 25th and
75th percentiles (box bounds); whiskers extend to the most extreme value
within 1.5 × the interquartile range of the nearest box bound, and all
plants are plotted as points. The statistics above the plots are
estimated marginal mean (EMM) contrasts of the BX effect within each
genotype (ΔEMM = BXplus − BXminus), computed from a linear mixed-effects
model with experiment as a random intercept to account for variation
between experiments. Each ΔEMM is a two-sided t-test on the estimated
marginal means (emmeans). ΔEMM is therefore expressed in the units of
the measurement, not in z-score units. (c) Summary of the same linear
mixed-effects model in (b). For each genotype, the BX effect is reported
as the EMM contrast Δ(BXplus − BXminus) with unadjusted two-sided P
values. The table also reports interaction contrasts relative to WT (ΔΔ
vs WT = Δ\_mutant − Δ\_WT), which test whether the BX response of each
mutant differs from that of WT; these three comparisons form one family
per measure and are two-sided t-tests with Holm adjusted P values. Data
are from Dataset 4.

\subsubsection{Plot combined ETI fresh
weight}\label{plot-combined-eti-fresh-weight}

\begin{figure}

\centering{

\pandocbounded{\includegraphics[keepaspectratio]{04_Mutant_screening_files/figure-pdf/fig-eti-fw-combined-1.pdf}}

}

\caption{\label{fig-eti-fw-combined}Extended Data Figure 6b (fresh
weight) \textbar{} Growth feedback in a BXplus microbiome remains intact
in effector-triggered immunity signalling mutants. (a) Phenotypes of
6-week-old WT, and mutants (eds1-12, pad4-1 and sag101-3) affected in
effector-triggered immunity signalling pathways grown on the BXminus and
BXplus soils. Scale bar, 2 cm. (b) Z-score--transformed fresh weight and
rosette area of WT, eds1-12, pad4-1, and sag101-3 plants across four
independent experiments illustrate the positive growth feedback on the
BXplus soil microbiome. Values were z-score--transformed within
experiment to account for between-experiment variation. Each point is an
individual plant grown in its own pot, i.e.~a biological replicate;
number of replicates is given below each plot. Boxplots show the median
(centre line) and the 25th and 75th percentiles (box bounds); whiskers
extend to the most extreme value within 1.5 × the interquartile range of
the nearest box bound, and all plants are plotted as points. The
statistics above the plots are estimated marginal mean (EMM) contrasts
of the BX effect within each genotype (ΔEMM = BXplus − BXminus),
computed from a linear mixed-effects model with experiment as a random
intercept to account for variation between experiments. Each ΔEMM is a
two-sided t-test on the estimated marginal means (emmeans). ΔEMM is
therefore expressed in the units of the measurement, not in z-score
units. (c) Summary of the same linear mixed-effects model in (b). For
each genotype, the BX effect is reported as the EMM contrast Δ(BXplus −
BXminus) with unadjusted two-sided P values. The table also reports
interaction contrasts relative to WT (ΔΔ vs WT = Δ\_mutant − Δ\_WT),
which test whether the BX response of each mutant differs from that of
WT; these three comparisons form one family per measure and are
two-sided t-tests with Holm adjusted P values. Data are from Dataset 4.}

\end{figure}%

\subsubsection{LMM and EMM: rosette
area}\label{lmm-and-emm-rosette-area}

\begin{verbatim}
Linear mixed model fit by REML. t-tests use Satterthwaite's method [
lmerModLmerTest]
Formula: formula_lmm
   Data: data

REML criterion at convergence: 11287.5

Scaled residuals: 
    Min      1Q  Median      3Q     Max 
-2.8279 -0.5525 -0.0352  0.5055  4.2185 

Random effects:
 Groups     Name        Variance  Std.Dev.
 experiment (Intercept) 7.732e+11 879334  
 Residual               1.426e+11 377571  
Number of obs: 402, groups:  experiment, 4

Fixed effects:
                    Estimate Std. Error         df t value Pr(>|t|)   
(Intercept)       990741.199 442776.214      3.075   2.238  0.10905   
BXplus            243686.909  73728.413    391.001   3.305  0.00104 **
lineeds1          113047.669  74075.920    391.001   1.526  0.12779   
linepad4           29914.915  77450.254    391.004   0.386  0.69952   
linesag101        -14491.448  75596.489    391.001  -0.192  0.84808   
BXplus:lineeds1      653.973 104509.216    391.001   0.006  0.99501   
BXplus:linepad4    -2083.262 107949.685    391.002  -0.019  0.98461   
BXplus:linesag101 -20285.286 105380.519    391.002  -0.192  0.84745   
---
Signif. codes:  0 '***' 0.001 '**' 0.01 '*' 0.05 '.' 0.1 ' ' 1

Correlation of Fixed Effects:
            (Intr) BXplus linds1 linpd4 lns101 BXpl:1 BXpl:4
BXplus      -0.084                                          
lineeds1    -0.084  0.502                                   
linepad4    -0.080  0.480  0.479                            
linesag101  -0.082  0.492  0.490  0.469                     
BXpls:lnds1  0.059 -0.706 -0.708 -0.339 -0.347              
BXpls:lnpd4  0.057 -0.683 -0.343 -0.716 -0.336  0.482       
BXpls:ln101  0.059 -0.700 -0.351 -0.336 -0.717  0.494  0.478
\end{verbatim}

\begin{longtable}[]{@{}
  >{\raggedright\arraybackslash}p{(\linewidth - 16\tabcolsep) * \real{0.1429}}
  >{\raggedright\arraybackslash}p{(\linewidth - 16\tabcolsep) * \real{0.0769}}
  >{\raggedleft\arraybackslash}p{(\linewidth - 16\tabcolsep) * \real{0.1978}}
  >{\raggedleft\arraybackslash}p{(\linewidth - 16\tabcolsep) * \real{0.0989}}
  >{\raggedleft\arraybackslash}p{(\linewidth - 16\tabcolsep) * \real{0.0989}}
  >{\raggedleft\arraybackslash}p{(\linewidth - 16\tabcolsep) * \real{0.0989}}
  >{\raggedleft\arraybackslash}p{(\linewidth - 16\tabcolsep) * \real{0.0989}}
  >{\raggedleft\arraybackslash}p{(\linewidth - 16\tabcolsep) * \real{0.0989}}
  >{\raggedleft\arraybackslash}p{(\linewidth - 16\tabcolsep) * \real{0.0879}}@{}}
\caption{Within-line ΔEMM = BXplus - BXminus for ETI rosette
area}\tabularnewline
\toprule\noalign{}
\begin{minipage}[b]{\linewidth}\raggedright
contrast
\end{minipage} & \begin{minipage}[b]{\linewidth}\raggedright
line
\end{minipage} & \begin{minipage}[b]{\linewidth}\raggedleft
effect\_plus\_minus
\end{minipage} & \begin{minipage}[b]{\linewidth}\raggedleft
SE
\end{minipage} & \begin{minipage}[b]{\linewidth}\raggedleft
df
\end{minipage} & \begin{minipage}[b]{\linewidth}\raggedleft
lower.CL
\end{minipage} & \begin{minipage}[b]{\linewidth}\raggedleft
upper.CL
\end{minipage} & \begin{minipage}[b]{\linewidth}\raggedleft
t
\end{minipage} & \begin{minipage}[b]{\linewidth}\raggedleft
p
\end{minipage} \\
\midrule\noalign{}
\endfirsthead
\toprule\noalign{}
\begin{minipage}[b]{\linewidth}\raggedright
contrast
\end{minipage} & \begin{minipage}[b]{\linewidth}\raggedright
line
\end{minipage} & \begin{minipage}[b]{\linewidth}\raggedleft
effect\_plus\_minus
\end{minipage} & \begin{minipage}[b]{\linewidth}\raggedleft
SE
\end{minipage} & \begin{minipage}[b]{\linewidth}\raggedleft
df
\end{minipage} & \begin{minipage}[b]{\linewidth}\raggedleft
lower.CL
\end{minipage} & \begin{minipage}[b]{\linewidth}\raggedleft
upper.CL
\end{minipage} & \begin{minipage}[b]{\linewidth}\raggedleft
t
\end{minipage} & \begin{minipage}[b]{\linewidth}\raggedleft
p
\end{minipage} \\
\midrule\noalign{}
\endhead
\bottomrule\noalign{}
\endlastfoot
plus - minus & WT & 243686.9 & 73728.49 & 391.0012 & 98733.03 & 388640.8
& 3.305193 & 0.00104 \\
plus - minus & eds1 & 244340.9 & 74062.12 & 391.0006 & 98731.07 &
389950.7 & 3.299134 & 0.00106 \\
plus - minus & pad4 & 241603.6 & 78827.63 & 391.0009 & 86624.61 &
396582.7 & 3.064961 & 0.00233 \\
plus - minus & sag101 & 223401.6 & 75249.23 & 391.0007 & 75457.89 &
371345.4 & 2.968822 & 0.00317 \\
\end{longtable}

Extended Data Figure 6c (rosette area) \textbar{} Growth feedback in a
BXplus microbiome remains intact in effector-triggered immunity
signalling mutants. (a) Phenotypes of 6-week-old WT, and mutants
(eds1-12, pad4-1 and sag101-3) affected in effector-triggered immunity
signalling pathways grown on the BXminus and BXplus soils. Scale bar, 2
cm. (b) Z-score--transformed fresh weight and rosette area of WT,
eds1-12, pad4-1, and sag101-3 plants across four independent experiments
illustrate the positive growth feedback on the BXplus soil microbiome.
Values were z-score--transformed within experiment to account for
between-experiment variation. Each point is an individual plant grown in
its own pot, i.e.~a biological replicate; number of replicates is given
below each plot. Boxplots show the median (centre line) and the 25th and
75th percentiles (box bounds); whiskers extend to the most extreme value
within 1.5 × the interquartile range of the nearest box bound, and all
plants are plotted as points. The statistics above the plots are
estimated marginal mean (EMM) contrasts of the BX effect within each
genotype (ΔEMM = BXplus − BXminus), computed from a linear mixed-effects
model with experiment as a random intercept to account for variation
between experiments. Each ΔEMM is a two-sided t-test on the estimated
marginal means (emmeans). ΔEMM is therefore expressed in the units of
the measurement, not in z-score units. (c) Summary of the same linear
mixed-effects model in (b). For each genotype, the BX effect is reported
as the EMM contrast Δ(BXplus − BXminus) with unadjusted two-sided P
values. The table also reports interaction contrasts relative to WT (ΔΔ
vs WT = Δ\_mutant − Δ\_WT), which test whether the BX response of each
mutant differs from that of WT; these three comparisons form one family
per measure and are two-sided t-tests with Holm adjusted P values. Data
are from Dataset 4.

\begin{longtable}[]{@{}
  >{\raggedright\arraybackslash}p{(\linewidth - 16\tabcolsep) * \real{0.1546}}
  >{\raggedright\arraybackslash}p{(\linewidth - 16\tabcolsep) * \real{0.1753}}
  >{\raggedleft\arraybackslash}p{(\linewidth - 16\tabcolsep) * \real{0.1134}}
  >{\raggedleft\arraybackslash}p{(\linewidth - 16\tabcolsep) * \real{0.0928}}
  >{\raggedleft\arraybackslash}p{(\linewidth - 16\tabcolsep) * \real{0.0928}}
  >{\raggedleft\arraybackslash}p{(\linewidth - 16\tabcolsep) * \real{0.1031}}
  >{\raggedleft\arraybackslash}p{(\linewidth - 16\tabcolsep) * \real{0.0928}}
  >{\raggedleft\arraybackslash}p{(\linewidth - 16\tabcolsep) * \real{0.1134}}
  >{\raggedleft\arraybackslash}p{(\linewidth - 16\tabcolsep) * \real{0.0619}}@{}}
\caption{Difference-in-differences relative to WT for ETI rosette
area}\tabularnewline
\toprule\noalign{}
\begin{minipage}[b]{\linewidth}\raggedright
BX\_revpairwise
\end{minipage} & \begin{minipage}[b]{\linewidth}\raggedright
line\_trt.vs.ctrl
\end{minipage} & \begin{minipage}[b]{\linewidth}\raggedleft
diff\_vs\_WT
\end{minipage} & \begin{minipage}[b]{\linewidth}\raggedleft
SE
\end{minipage} & \begin{minipage}[b]{\linewidth}\raggedleft
df
\end{minipage} & \begin{minipage}[b]{\linewidth}\raggedleft
lower.CL
\end{minipage} & \begin{minipage}[b]{\linewidth}\raggedleft
upper.CL
\end{minipage} & \begin{minipage}[b]{\linewidth}\raggedleft
t
\end{minipage} & \begin{minipage}[b]{\linewidth}\raggedleft
p\_adj
\end{minipage} \\
\midrule\noalign{}
\endfirsthead
\toprule\noalign{}
\begin{minipage}[b]{\linewidth}\raggedright
BX\_revpairwise
\end{minipage} & \begin{minipage}[b]{\linewidth}\raggedright
line\_trt.vs.ctrl
\end{minipage} & \begin{minipage}[b]{\linewidth}\raggedleft
diff\_vs\_WT
\end{minipage} & \begin{minipage}[b]{\linewidth}\raggedleft
SE
\end{minipage} & \begin{minipage}[b]{\linewidth}\raggedleft
df
\end{minipage} & \begin{minipage}[b]{\linewidth}\raggedleft
lower.CL
\end{minipage} & \begin{minipage}[b]{\linewidth}\raggedleft
upper.CL
\end{minipage} & \begin{minipage}[b]{\linewidth}\raggedleft
t
\end{minipage} & \begin{minipage}[b]{\linewidth}\raggedleft
p\_adj
\end{minipage} \\
\midrule\noalign{}
\endhead
\bottomrule\noalign{}
\endlastfoot
plus - minus & eds1 - WT & 653.973 & 104509.3 & 391.0011 & -250620.4 &
251928.4 & 0.0062576 & 1 \\
plus - minus & pad4 - WT & -2083.262 & 107949.8 & 391.0015 & -261629.8 &
257463.2 & -0.0192984 & 1 \\
plus - minus & sag101 - WT & -20285.286 & 105380.7 & 391.0019 &
-273654.7 & 233084.2 & -0.1924953 & 1 \\
\end{longtable}

Extended Data Figure 6c (rosette area) \textbar{} Growth feedback in a
BXplus microbiome remains intact in effector-triggered immunity
signalling mutants. (a) Phenotypes of 6-week-old WT, and mutants
(eds1-12, pad4-1 and sag101-3) affected in effector-triggered immunity
signalling pathways grown on the BXminus and BXplus soils. Scale bar, 2
cm. (b) Z-score--transformed fresh weight and rosette area of WT,
eds1-12, pad4-1, and sag101-3 plants across four independent experiments
illustrate the positive growth feedback on the BXplus soil microbiome.
Values were z-score--transformed within experiment to account for
between-experiment variation. Each point is an individual plant grown in
its own pot, i.e.~a biological replicate; number of replicates is given
below each plot. Boxplots show the median (centre line) and the 25th and
75th percentiles (box bounds); whiskers extend to the most extreme value
within 1.5 × the interquartile range of the nearest box bound, and all
plants are plotted as points. The statistics above the plots are
estimated marginal mean (EMM) contrasts of the BX effect within each
genotype (ΔEMM = BXplus − BXminus), computed from a linear mixed-effects
model with experiment as a random intercept to account for variation
between experiments. Each ΔEMM is a two-sided t-test on the estimated
marginal means (emmeans). ΔEMM is therefore expressed in the units of
the measurement, not in z-score units. (c) Summary of the same linear
mixed-effects model in (b). For each genotype, the BX effect is reported
as the EMM contrast Δ(BXplus − BXminus) with unadjusted two-sided P
values. The table also reports interaction contrasts relative to WT (ΔΔ
vs WT = Δ\_mutant − Δ\_WT), which test whether the BX response of each
mutant differs from that of WT; these three comparisons form one family
per measure and are two-sided t-tests with Holm adjusted P values. Data
are from Dataset 4.

\begin{longtable}[]{@{}
  >{\raggedright\arraybackslash}p{(\linewidth - 16\tabcolsep) * \real{0.1412}}
  >{\raggedright\arraybackslash}p{(\linewidth - 16\tabcolsep) * \real{0.0706}}
  >{\raggedleft\arraybackslash}p{(\linewidth - 16\tabcolsep) * \real{0.1176}}
  >{\raggedleft\arraybackslash}p{(\linewidth - 16\tabcolsep) * \real{0.1059}}
  >{\raggedleft\arraybackslash}p{(\linewidth - 16\tabcolsep) * \real{0.1059}}
  >{\raggedleft\arraybackslash}p{(\linewidth - 16\tabcolsep) * \real{0.1294}}
  >{\raggedleft\arraybackslash}p{(\linewidth - 16\tabcolsep) * \real{0.1059}}
  >{\raggedleft\arraybackslash}p{(\linewidth - 16\tabcolsep) * \real{0.1294}}
  >{\raggedleft\arraybackslash}p{(\linewidth - 16\tabcolsep) * \real{0.0941}}@{}}
\caption{Mutant vs WT contrasts within each BX condition for ETI rosette
area (Holm-adjusted)}\tabularnewline
\toprule\noalign{}
\begin{minipage}[b]{\linewidth}\raggedright
contrast
\end{minipage} & \begin{minipage}[b]{\linewidth}\raggedright
BX
\end{minipage} & \begin{minipage}[b]{\linewidth}\raggedleft
estimate
\end{minipage} & \begin{minipage}[b]{\linewidth}\raggedleft
SE
\end{minipage} & \begin{minipage}[b]{\linewidth}\raggedleft
df
\end{minipage} & \begin{minipage}[b]{\linewidth}\raggedleft
lower.CL
\end{minipage} & \begin{minipage}[b]{\linewidth}\raggedleft
upper.CL
\end{minipage} & \begin{minipage}[b]{\linewidth}\raggedleft
t.ratio
\end{minipage} & \begin{minipage}[b]{\linewidth}\raggedleft
p.value
\end{minipage} \\
\midrule\noalign{}
\endfirsthead
\toprule\noalign{}
\begin{minipage}[b]{\linewidth}\raggedright
contrast
\end{minipage} & \begin{minipage}[b]{\linewidth}\raggedright
BX
\end{minipage} & \begin{minipage}[b]{\linewidth}\raggedleft
estimate
\end{minipage} & \begin{minipage}[b]{\linewidth}\raggedleft
SE
\end{minipage} & \begin{minipage}[b]{\linewidth}\raggedleft
df
\end{minipage} & \begin{minipage}[b]{\linewidth}\raggedleft
lower.CL
\end{minipage} & \begin{minipage}[b]{\linewidth}\raggedleft
upper.CL
\end{minipage} & \begin{minipage}[b]{\linewidth}\raggedleft
t.ratio
\end{minipage} & \begin{minipage}[b]{\linewidth}\raggedleft
p.value
\end{minipage} \\
\midrule\noalign{}
\endhead
\bottomrule\noalign{}
\endlastfoot
eds1 - WT & minus & 113047.67 & 74075.99 & 391.0011 & -65055.11 &
291150.4 & 1.5261041 & 0.383 \\
pad4 - WT & minus & 29914.92 & 77450.52 & 391.0040 & -156301.33 &
216131.2 & 0.3862455 & 1.000 \\
sag101 - WT & minus & -14491.45 & 75596.54 & 391.0008 & -196250.13 &
167267.2 & -0.1916946 & 1.000 \\
eds1 - WT & plus & 113701.64 & 73764.74 & 391.0027 & -63652.78 &
291056.1 & 1.5414092 & 0.372 \\
pad4 - WT & plus & 27831.65 & 75400.74 & 391.0065 & -153456.24 &
209119.5 & 0.3691164 & 1.000 \\
sag101 - WT & plus & -34776.73 & 73482.43 & 391.0055 & -211452.40 &
141898.9 & -0.4732660 & 1.000 \\
\end{longtable}

Extended Data Figure 6c (rosette area) \textbar{} Growth feedback in a
BXplus microbiome remains intact in effector-triggered immunity
signalling mutants. (a) Phenotypes of 6-week-old WT, and mutants
(eds1-12, pad4-1 and sag101-3) affected in effector-triggered immunity
signalling pathways grown on the BXminus and BXplus soils. Scale bar, 2
cm. (b) Z-score--transformed fresh weight and rosette area of WT,
eds1-12, pad4-1, and sag101-3 plants across four independent experiments
illustrate the positive growth feedback on the BXplus soil microbiome.
Values were z-score--transformed within experiment to account for
between-experiment variation. Each point is an individual plant grown in
its own pot, i.e.~a biological replicate; number of replicates is given
below each plot. Boxplots show the median (centre line) and the 25th and
75th percentiles (box bounds); whiskers extend to the most extreme value
within 1.5 × the interquartile range of the nearest box bound, and all
plants are plotted as points. The statistics above the plots are
estimated marginal mean (EMM) contrasts of the BX effect within each
genotype (ΔEMM = BXplus − BXminus), computed from a linear mixed-effects
model with experiment as a random intercept to account for variation
between experiments. Each ΔEMM is a two-sided t-test on the estimated
marginal means (emmeans). ΔEMM is therefore expressed in the units of
the measurement, not in z-score units. (c) Summary of the same linear
mixed-effects model in (b). For each genotype, the BX effect is reported
as the EMM contrast Δ(BXplus − BXminus) with unadjusted two-sided P
values. The table also reports interaction contrasts relative to WT (ΔΔ
vs WT = Δ\_mutant − Δ\_WT), which test whether the BX response of each
mutant differs from that of WT; these three comparisons form one family
per measure and are two-sided t-tests with Holm adjusted P values. Data
are from Dataset 4.

\subsubsection{Plot combined ETI rosette
area}\label{plot-combined-eti-rosette-area}

\begin{figure}

\centering{

\pandocbounded{\includegraphics[keepaspectratio]{04_Mutant_screening_files/figure-pdf/fig-eti-area-combined-1.pdf}}

}

\caption{\label{fig-eti-area-combined}Extended Data Figure 6b (rosette
area) \textbar{} Growth feedback in a BXplus microbiome remains intact
in effector-triggered immunity signalling mutants. (a) Phenotypes of
6-week-old WT, and mutants (eds1-12, pad4-1 and sag101-3) affected in
effector-triggered immunity signalling pathways grown on the BXminus and
BXplus soils. Scale bar, 2 cm. (b) Z-score--transformed fresh weight and
rosette area of WT, eds1-12, pad4-1, and sag101-3 plants across four
independent experiments illustrate the positive growth feedback on the
BXplus soil microbiome. Values were z-score--transformed within
experiment to account for between-experiment variation. Each point is an
individual plant grown in its own pot, i.e.~a biological replicate;
number of replicates is given below each plot. Boxplots show the median
(centre line) and the 25th and 75th percentiles (box bounds); whiskers
extend to the most extreme value within 1.5 × the interquartile range of
the nearest box bound, and all plants are plotted as points. The
statistics above the plots are estimated marginal mean (EMM) contrasts
of the BX effect within each genotype (ΔEMM = BXplus − BXminus),
computed from a linear mixed-effects model with experiment as a random
intercept to account for variation between experiments. Each ΔEMM is a
two-sided t-test on the estimated marginal means (emmeans). ΔEMM is
therefore expressed in the units of the measurement, not in z-score
units. (c) Summary of the same linear mixed-effects model in (b). For
each genotype, the BX effect is reported as the EMM contrast Δ(BXplus −
BXminus) with unadjusted two-sided P values. The table also reports
interaction contrasts relative to WT (ΔΔ vs WT = Δ\_mutant − Δ\_WT),
which test whether the BX response of each mutant differs from that of
WT; these three comparisons form one family per measure and are
two-sided t-tests with Holm adjusted P values. Data are from Dataset 4.}

\end{figure}%

\newpage

\section{\texorpdfstring{\textbf{Combined mmf1 analysis across E-III,
E-IV, E-V, E-VIII and
E-IX}}{Combined mmf1 analysis across E-III, E-IV, E-V, E-VIII and E-IX}}\label{combined-mmf1-analysis-across-e-iii-e-iv-e-v-e-viii-and-e-ix}

\subsubsection{\texorpdfstring{\textbf{LMM and EMM: fresh
weight}}{LMM and EMM: fresh weight}}\label{lmm-and-emm-fresh-weight-1}

\begin{verbatim}
Linear mixed model fit by REML. t-tests use Satterthwaite's method [
lmerModLmerTest]
Formula: formula_lmm
   Data: data

REML criterion at convergence: -767.4

Scaled residuals: 
    Min      1Q  Median      3Q     Max 
-4.7563 -0.4404  0.0308  0.4323  4.4196 

Random effects:
 Groups     Name        Variance Std.Dev.
 experiment (Intercept) 0.045798 0.2140  
 Residual               0.006691 0.0818  
Number of obs: 385, groups:  experiment, 6

Fixed effects:
                   Estimate Std. Error        df t value Pr(>|t|)    
(Intercept)         0.26500    0.08783   5.07779   3.017 0.028943 *  
BXplus              0.05276    0.01263 373.99689   4.178 3.66e-05 ***
linemmf1-1         -0.05736    0.01659 374.18448  -3.457 0.000609 ***
linemmf1-2         -0.02057    0.01355 374.06135  -1.518 0.129884    
BXplus:linemmf1-1  -0.04445    0.02242 373.99666  -1.982 0.048170 *  
BXplus:linemmf1-2  -0.05779    0.01877 373.99763  -3.079 0.002230 ** 
---
Signif. codes:  0 '***' 0.001 '**' 0.01 '*' 0.05 '.' 0.1 ' ' 1

Correlation of Fixed Effects:
            (Intr) BXplus lnm1-1 lnm1-2 BX:1-1
BXplus      -0.074                            
linemmf1-1  -0.055  0.385                     
linemmf1-2  -0.069  0.477  0.330              
BXpls:ln1-1  0.041 -0.563 -0.677 -0.271       
BXpls:ln1-2  0.049 -0.673 -0.260 -0.702  0.379
\end{verbatim}

\begin{longtable}[]{@{}
  >{\raggedright\arraybackslash}p{(\linewidth - 16\tabcolsep) * \real{0.1429}}
  >{\raggedright\arraybackslash}p{(\linewidth - 16\tabcolsep) * \real{0.0769}}
  >{\raggedleft\arraybackslash}p{(\linewidth - 16\tabcolsep) * \real{0.1978}}
  >{\raggedleft\arraybackslash}p{(\linewidth - 16\tabcolsep) * \real{0.0659}}
  >{\raggedleft\arraybackslash}p{(\linewidth - 16\tabcolsep) * \real{0.0989}}
  >{\raggedleft\arraybackslash}p{(\linewidth - 16\tabcolsep) * \real{0.0989}}
  >{\raggedleft\arraybackslash}p{(\linewidth - 16\tabcolsep) * \real{0.0989}}
  >{\raggedleft\arraybackslash}p{(\linewidth - 16\tabcolsep) * \real{0.1209}}
  >{\raggedleft\arraybackslash}p{(\linewidth - 16\tabcolsep) * \real{0.0989}}@{}}
\caption{Within-line ΔEMM = BXplus - BXminus for mmf1 fresh
weight}\tabularnewline
\toprule\noalign{}
\begin{minipage}[b]{\linewidth}\raggedright
contrast
\end{minipage} & \begin{minipage}[b]{\linewidth}\raggedright
line
\end{minipage} & \begin{minipage}[b]{\linewidth}\raggedleft
effect\_plus\_minus
\end{minipage} & \begin{minipage}[b]{\linewidth}\raggedleft
SE
\end{minipage} & \begin{minipage}[b]{\linewidth}\raggedleft
df
\end{minipage} & \begin{minipage}[b]{\linewidth}\raggedleft
lower.CL
\end{minipage} & \begin{minipage}[b]{\linewidth}\raggedleft
upper.CL
\end{minipage} & \begin{minipage}[b]{\linewidth}\raggedleft
t
\end{minipage} & \begin{minipage}[b]{\linewidth}\raggedleft
p
\end{minipage} \\
\midrule\noalign{}
\endfirsthead
\toprule\noalign{}
\begin{minipage}[b]{\linewidth}\raggedright
contrast
\end{minipage} & \begin{minipage}[b]{\linewidth}\raggedright
line
\end{minipage} & \begin{minipage}[b]{\linewidth}\raggedleft
effect\_plus\_minus
\end{minipage} & \begin{minipage}[b]{\linewidth}\raggedleft
SE
\end{minipage} & \begin{minipage}[b]{\linewidth}\raggedleft
df
\end{minipage} & \begin{minipage}[b]{\linewidth}\raggedleft
lower.CL
\end{minipage} & \begin{minipage}[b]{\linewidth}\raggedleft
upper.CL
\end{minipage} & \begin{minipage}[b]{\linewidth}\raggedleft
t
\end{minipage} & \begin{minipage}[b]{\linewidth}\raggedleft
p
\end{minipage} \\
\midrule\noalign{}
\endhead
\bottomrule\noalign{}
\endlastfoot
plus - minus & WT & 0.053 & 0.013 & 374.0009 & 0.028 & 0.078 & 4.1781605
& 3.66e-05 \\
plus - minus & mmf1-1 & 0.008 & 0.019 & 374.0007 & -0.028 & 0.045 &
0.4489470 & 6.54e-01 \\
plus - minus & mmf1-2 & -0.005 & 0.014 & 374.0026 & -0.032 & 0.022 &
-0.3619978 & 7.18e-01 \\
\end{longtable}

Extended Data Figure 5b (fresh weight) \textbar{} Growth feedback in a
BXplus microbiome is consistently lost in mmf1 mutants. (a and b)
Repeated growth feedbacks (Experiment III - V, VIII, IX and X) of Col-0
wild type (WT) and two mmf1 T-DNA insertion mutants (At3g44630) grown on
the BXminus and BXplus soil microbiomes. Feedback experiments were
performed on greenhouse-conditioned soil batches BS04 (Experiment III
and IV), BS09 (V and X), BS10 (IX) and BS11 (VIII). (a)
Z-score--transformed fresh weight and rosette area of WT and mmf1
mutants across six independent experiments. WT was present in all six
experiments, mmf1-1 in three and mmf1-2 in five. Values were
z-score--transformed within experiment to account for between-experiment
variation. Each point is an individual plant grown in its own pot,
i.e.~a biological replicate; number of replicates is given below each
plot. Boxplots show the median (centre line) and the 25th and 75th
percentiles (box bounds); whiskers extend to the most extreme value
within 1.5 × the interquartile range of the nearest box bound, and all
plants are plotted as points. The statistics above the plots are
estimated marginal mean (EMM) contrasts of the BX effect within each
genotype (ΔEMM = BXplus − BXminus), computed from a linear mixed-effects
model with experiment as a random intercept to account for variation
between experiments. Each ΔEMM is a two-sided t-test on the estimated
marginal means (emmeans). ΔEMM is therefore expressed in the units of
the measurement, not in z-score units. (b) Summary of the same linear
mixed-effects model in (a). For each genotype, the BX effect is reported
as the EMM contrast Δ(BXplus − BXminus) with unadjusted two-sided P
values. The table also reports interaction contrasts relative to WT (ΔΔ
vs WT = Δ\_mutant − Δ\_WT), which test whether the BX response of each
mutant differs from that of WT; these two comparisons form one family
per measure and are two-sided t-tests with Holm adjusted P values. Data
are from Dataset 4.

\begin{longtable}[]{@{}
  >{\raggedright\arraybackslash}p{(\linewidth - 16\tabcolsep) * \real{0.1596}}
  >{\raggedright\arraybackslash}p{(\linewidth - 16\tabcolsep) * \real{0.1809}}
  >{\raggedleft\arraybackslash}p{(\linewidth - 16\tabcolsep) * \real{0.1170}}
  >{\raggedleft\arraybackslash}p{(\linewidth - 16\tabcolsep) * \real{0.0638}}
  >{\raggedleft\arraybackslash}p{(\linewidth - 16\tabcolsep) * \real{0.0957}}
  >{\raggedleft\arraybackslash}p{(\linewidth - 16\tabcolsep) * \real{0.0957}}
  >{\raggedleft\arraybackslash}p{(\linewidth - 16\tabcolsep) * \real{0.0957}}
  >{\raggedleft\arraybackslash}p{(\linewidth - 16\tabcolsep) * \real{0.1064}}
  >{\raggedleft\arraybackslash}p{(\linewidth - 16\tabcolsep) * \real{0.0851}}@{}}
\caption{Difference-in-differences relative to WT for mmf1 fresh
weight}\tabularnewline
\toprule\noalign{}
\begin{minipage}[b]{\linewidth}\raggedright
BX\_revpairwise
\end{minipage} & \begin{minipage}[b]{\linewidth}\raggedright
line\_trt.vs.ctrl
\end{minipage} & \begin{minipage}[b]{\linewidth}\raggedleft
diff\_vs\_WT
\end{minipage} & \begin{minipage}[b]{\linewidth}\raggedleft
SE
\end{minipage} & \begin{minipage}[b]{\linewidth}\raggedleft
df
\end{minipage} & \begin{minipage}[b]{\linewidth}\raggedleft
lower.CL
\end{minipage} & \begin{minipage}[b]{\linewidth}\raggedleft
upper.CL
\end{minipage} & \begin{minipage}[b]{\linewidth}\raggedleft
t
\end{minipage} & \begin{minipage}[b]{\linewidth}\raggedleft
p\_adj
\end{minipage} \\
\midrule\noalign{}
\endfirsthead
\toprule\noalign{}
\begin{minipage}[b]{\linewidth}\raggedright
BX\_revpairwise
\end{minipage} & \begin{minipage}[b]{\linewidth}\raggedright
line\_trt.vs.ctrl
\end{minipage} & \begin{minipage}[b]{\linewidth}\raggedleft
diff\_vs\_WT
\end{minipage} & \begin{minipage}[b]{\linewidth}\raggedleft
SE
\end{minipage} & \begin{minipage}[b]{\linewidth}\raggedleft
df
\end{minipage} & \begin{minipage}[b]{\linewidth}\raggedleft
lower.CL
\end{minipage} & \begin{minipage}[b]{\linewidth}\raggedleft
upper.CL
\end{minipage} & \begin{minipage}[b]{\linewidth}\raggedleft
t
\end{minipage} & \begin{minipage}[b]{\linewidth}\raggedleft
p\_adj
\end{minipage} \\
\midrule\noalign{}
\endhead
\bottomrule\noalign{}
\endlastfoot
plus - minus & (mmf1-1) - WT & -0.044 & 0.022 & 374.0007 & -0.095 &
0.006 & -1.982364 & 0.04820 \\
plus - minus & (mmf1-2) - WT & -0.058 & 0.019 & 374.0016 & -0.100 &
-0.016 & -3.079076 & 0.00446 \\
\end{longtable}

Extended Data Figure 5b (fresh weight) \textbar{} Growth feedback in a
BXplus microbiome is consistently lost in mmf1 mutants. (a and b)
Repeated growth feedbacks (Experiment III - V, VIII, IX and X) of Col-0
wild type (WT) and two mmf1 T-DNA insertion mutants (At3g44630) grown on
the BXminus and BXplus soil microbiomes. Feedback experiments were
performed on greenhouse-conditioned soil batches BS04 (Experiment III
and IV), BS09 (V and X), BS10 (IX) and BS11 (VIII). (a)
Z-score--transformed fresh weight and rosette area of WT and mmf1
mutants across six independent experiments. WT was present in all six
experiments, mmf1-1 in three and mmf1-2 in five. Values were
z-score--transformed within experiment to account for between-experiment
variation. Each point is an individual plant grown in its own pot,
i.e.~a biological replicate; number of replicates is given below each
plot. Boxplots show the median (centre line) and the 25th and 75th
percentiles (box bounds); whiskers extend to the most extreme value
within 1.5 × the interquartile range of the nearest box bound, and all
plants are plotted as points. The statistics above the plots are
estimated marginal mean (EMM) contrasts of the BX effect within each
genotype (ΔEMM = BXplus − BXminus), computed from a linear mixed-effects
model with experiment as a random intercept to account for variation
between experiments. Each ΔEMM is a two-sided t-test on the estimated
marginal means (emmeans). ΔEMM is therefore expressed in the units of
the measurement, not in z-score units. (b) Summary of the same linear
mixed-effects model in (a). For each genotype, the BX effect is reported
as the EMM contrast Δ(BXplus − BXminus) with unadjusted two-sided P
values. The table also reports interaction contrasts relative to WT (ΔΔ
vs WT = Δ\_mutant − Δ\_WT), which test whether the BX response of each
mutant differs from that of WT; these two comparisons form one family
per measure and are two-sided t-tests with Holm adjusted P values. Data
are from Dataset 4.

\begin{longtable}[]{@{}
  >{\raggedright\arraybackslash}p{(\linewidth - 16\tabcolsep) * \real{0.1750}}
  >{\raggedright\arraybackslash}p{(\linewidth - 16\tabcolsep) * \real{0.0750}}
  >{\raggedleft\arraybackslash}p{(\linewidth - 16\tabcolsep) * \real{0.1125}}
  >{\raggedleft\arraybackslash}p{(\linewidth - 16\tabcolsep) * \real{0.0750}}
  >{\raggedleft\arraybackslash}p{(\linewidth - 16\tabcolsep) * \real{0.1125}}
  >{\raggedleft\arraybackslash}p{(\linewidth - 16\tabcolsep) * \real{0.1125}}
  >{\raggedleft\arraybackslash}p{(\linewidth - 16\tabcolsep) * \real{0.1125}}
  >{\raggedleft\arraybackslash}p{(\linewidth - 16\tabcolsep) * \real{0.1250}}
  >{\raggedleft\arraybackslash}p{(\linewidth - 16\tabcolsep) * \real{0.1000}}@{}}
\caption{Mutant vs WT contrasts within each BX condition for mmf1 fresh
weight (Holm-adjusted)}\tabularnewline
\toprule\noalign{}
\begin{minipage}[b]{\linewidth}\raggedright
contrast
\end{minipage} & \begin{minipage}[b]{\linewidth}\raggedright
BX
\end{minipage} & \begin{minipage}[b]{\linewidth}\raggedleft
estimate
\end{minipage} & \begin{minipage}[b]{\linewidth}\raggedleft
SE
\end{minipage} & \begin{minipage}[b]{\linewidth}\raggedleft
df
\end{minipage} & \begin{minipage}[b]{\linewidth}\raggedleft
lower.CL
\end{minipage} & \begin{minipage}[b]{\linewidth}\raggedleft
upper.CL
\end{minipage} & \begin{minipage}[b]{\linewidth}\raggedleft
t.ratio
\end{minipage} & \begin{minipage}[b]{\linewidth}\raggedleft
p.value
\end{minipage} \\
\midrule\noalign{}
\endfirsthead
\toprule\noalign{}
\begin{minipage}[b]{\linewidth}\raggedright
contrast
\end{minipage} & \begin{minipage}[b]{\linewidth}\raggedright
BX
\end{minipage} & \begin{minipage}[b]{\linewidth}\raggedleft
estimate
\end{minipage} & \begin{minipage}[b]{\linewidth}\raggedleft
SE
\end{minipage} & \begin{minipage}[b]{\linewidth}\raggedleft
df
\end{minipage} & \begin{minipage}[b]{\linewidth}\raggedleft
lower.CL
\end{minipage} & \begin{minipage}[b]{\linewidth}\raggedleft
upper.CL
\end{minipage} & \begin{minipage}[b]{\linewidth}\raggedleft
t.ratio
\end{minipage} & \begin{minipage}[b]{\linewidth}\raggedleft
p.value
\end{minipage} \\
\midrule\noalign{}
\endhead
\bottomrule\noalign{}
\endlastfoot
(mmf1-1) - WT & minus & -0.057 & 0.017 & 374.1883 & -0.095 & -0.020 &
-3.456509 & 0.00122 \\
(mmf1-2) - WT & minus & -0.021 & 0.014 & 374.0653 & -0.051 & 0.010 &
-1.517848 & 0.13000 \\
(mmf1-1) - WT & plus & -0.102 & 0.017 & 374.2129 & -0.139 & -0.065 &
-6.149429 & 0.00000 \\
(mmf1-2) - WT & plus & -0.078 & 0.013 & 374.0675 & -0.108 & -0.048 &
-5.863093 & 0.00000 \\
\end{longtable}

Extended Data Figure 5b (fresh weight) \textbar{} Growth feedback in a
BXplus microbiome is consistently lost in mmf1 mutants. (a and b)
Repeated growth feedbacks (Experiment III - V, VIII, IX and X) of Col-0
wild type (WT) and two mmf1 T-DNA insertion mutants (At3g44630) grown on
the BXminus and BXplus soil microbiomes. Feedback experiments were
performed on greenhouse-conditioned soil batches BS04 (Experiment III
and IV), BS09 (V and X), BS10 (IX) and BS11 (VIII). (a)
Z-score--transformed fresh weight and rosette area of WT and mmf1
mutants across six independent experiments. WT was present in all six
experiments, mmf1-1 in three and mmf1-2 in five. Values were
z-score--transformed within experiment to account for between-experiment
variation. Each point is an individual plant grown in its own pot,
i.e.~a biological replicate; number of replicates is given below each
plot. Boxplots show the median (centre line) and the 25th and 75th
percentiles (box bounds); whiskers extend to the most extreme value
within 1.5 × the interquartile range of the nearest box bound, and all
plants are plotted as points. The statistics above the plots are
estimated marginal mean (EMM) contrasts of the BX effect within each
genotype (ΔEMM = BXplus − BXminus), computed from a linear mixed-effects
model with experiment as a random intercept to account for variation
between experiments. Each ΔEMM is a two-sided t-test on the estimated
marginal means (emmeans). ΔEMM is therefore expressed in the units of
the measurement, not in z-score units. (b) Summary of the same linear
mixed-effects model in (a). For each genotype, the BX effect is reported
as the EMM contrast Δ(BXplus − BXminus) with unadjusted two-sided P
values. The table also reports interaction contrasts relative to WT (ΔΔ
vs WT = Δ\_mutant − Δ\_WT), which test whether the BX response of each
mutant differs from that of WT; these two comparisons form one family
per measure and are two-sided t-tests with Holm adjusted P values. Data
are from Dataset 4.

\begin{longtable}[]{@{}lrrrrr@{}}
\caption{Estimated marginal means for mmf1 fresh weight by line,
averaged across BX conditions (experiment included as a random
intercept)}\tabularnewline
\toprule\noalign{}
line & emmean & SE & df & lower.CL & upper.CL \\
\midrule\noalign{}
\endfirsthead
\toprule\noalign{}
line & emmean & SE & df & lower.CL & upper.CL \\
\midrule\noalign{}
\endhead
\bottomrule\noalign{}
\endlastfoot
WT & 0.291 & 0.088 & 5.027000 & 0.067 & 0.516 \\
mmf1-1 & 0.212 & 0.088 & 5.119780 & -0.013 & 0.436 \\
mmf1-2 & 0.242 & 0.088 & 5.043722 & 0.017 & 0.467 \\
\end{longtable}

Extended Data Figure 5b (fresh weight) \textbar{} Growth feedback in a
BXplus microbiome is consistently lost in mmf1 mutants. (a and b)
Repeated growth feedbacks (Experiment III - V, VIII, IX and X) of Col-0
wild type (WT) and two mmf1 T-DNA insertion mutants (At3g44630) grown on
the BXminus and BXplus soil microbiomes. Feedback experiments were
performed on greenhouse-conditioned soil batches BS04 (Experiment III
and IV), BS09 (V and X), BS10 (IX) and BS11 (VIII). (a)
Z-score--transformed fresh weight and rosette area of WT and mmf1
mutants across six independent experiments. WT was present in all six
experiments, mmf1-1 in three and mmf1-2 in five. Values were
z-score--transformed within experiment to account for between-experiment
variation. Each point is an individual plant grown in its own pot,
i.e.~a biological replicate; number of replicates is given below each
plot. Boxplots show the median (centre line) and the 25th and 75th
percentiles (box bounds); whiskers extend to the most extreme value
within 1.5 × the interquartile range of the nearest box bound, and all
plants are plotted as points. The statistics above the plots are
estimated marginal mean (EMM) contrasts of the BX effect within each
genotype (ΔEMM = BXplus − BXminus), computed from a linear mixed-effects
model with experiment as a random intercept to account for variation
between experiments. Each ΔEMM is a two-sided t-test on the estimated
marginal means (emmeans). ΔEMM is therefore expressed in the units of
the measurement, not in z-score units. (b) Summary of the same linear
mixed-effects model in (a). For each genotype, the BX effect is reported
as the EMM contrast Δ(BXplus − BXminus) with unadjusted two-sided P
values. The table also reports interaction contrasts relative to WT (ΔΔ
vs WT = Δ\_mutant − Δ\_WT), which test whether the BX response of each
mutant differs from that of WT; these two comparisons form one family
per measure and are two-sided t-tests with Holm adjusted P values. Data
are from Dataset 4.

\begin{longtable}[]{@{}
  >{\raggedright\arraybackslash}p{(\linewidth - 14\tabcolsep) * \real{0.1944}}
  >{\raggedleft\arraybackslash}p{(\linewidth - 14\tabcolsep) * \real{0.1250}}
  >{\raggedleft\arraybackslash}p{(\linewidth - 14\tabcolsep) * \real{0.0833}}
  >{\raggedleft\arraybackslash}p{(\linewidth - 14\tabcolsep) * \real{0.1250}}
  >{\raggedleft\arraybackslash}p{(\linewidth - 14\tabcolsep) * \real{0.1250}}
  >{\raggedleft\arraybackslash}p{(\linewidth - 14\tabcolsep) * \real{0.1250}}
  >{\raggedleft\arraybackslash}p{(\linewidth - 14\tabcolsep) * \real{0.1389}}
  >{\raggedleft\arraybackslash}p{(\linewidth - 14\tabcolsep) * \real{0.0833}}@{}}
\caption{Mutant vs WT contrasts for mmf1 fresh weight, averaged across
BX conditions and corrected for experiment
(Holm-adjusted)}\tabularnewline
\toprule\noalign{}
\begin{minipage}[b]{\linewidth}\raggedright
contrast
\end{minipage} & \begin{minipage}[b]{\linewidth}\raggedleft
diff\_EMM
\end{minipage} & \begin{minipage}[b]{\linewidth}\raggedleft
SE
\end{minipage} & \begin{minipage}[b]{\linewidth}\raggedleft
df
\end{minipage} & \begin{minipage}[b]{\linewidth}\raggedleft
lower.CL
\end{minipage} & \begin{minipage}[b]{\linewidth}\raggedleft
upper.CL
\end{minipage} & \begin{minipage}[b]{\linewidth}\raggedleft
t
\end{minipage} & \begin{minipage}[b]{\linewidth}\raggedleft
p
\end{minipage} \\
\midrule\noalign{}
\endfirsthead
\toprule\noalign{}
\begin{minipage}[b]{\linewidth}\raggedright
contrast
\end{minipage} & \begin{minipage}[b]{\linewidth}\raggedleft
diff\_EMM
\end{minipage} & \begin{minipage}[b]{\linewidth}\raggedleft
SE
\end{minipage} & \begin{minipage}[b]{\linewidth}\raggedleft
df
\end{minipage} & \begin{minipage}[b]{\linewidth}\raggedleft
lower.CL
\end{minipage} & \begin{minipage}[b]{\linewidth}\raggedleft
upper.CL
\end{minipage} & \begin{minipage}[b]{\linewidth}\raggedleft
t
\end{minipage} & \begin{minipage}[b]{\linewidth}\raggedleft
p
\end{minipage} \\
\midrule\noalign{}
\endhead
\bottomrule\noalign{}
\endlastfoot
(mmf1-1) - WT & -0.080 & 0.012 & 374.3661 & -0.107 & -0.052 & -6.518329
& 0e+00 \\
(mmf1-2) - WT & -0.049 & 0.010 & 374.1273 & -0.071 & -0.028 & -5.127185
& 5e-07 \\
\end{longtable}

Extended Data Figure 5b (fresh weight) \textbar{} Growth feedback in a
BXplus microbiome is consistently lost in mmf1 mutants. (a and b)
Repeated growth feedbacks (Experiment III - V, VIII, IX and X) of Col-0
wild type (WT) and two mmf1 T-DNA insertion mutants (At3g44630) grown on
the BXminus and BXplus soil microbiomes. Feedback experiments were
performed on greenhouse-conditioned soil batches BS04 (Experiment III
and IV), BS09 (V and X), BS10 (IX) and BS11 (VIII). (a)
Z-score--transformed fresh weight and rosette area of WT and mmf1
mutants across six independent experiments. WT was present in all six
experiments, mmf1-1 in three and mmf1-2 in five. Values were
z-score--transformed within experiment to account for between-experiment
variation. Each point is an individual plant grown in its own pot,
i.e.~a biological replicate; number of replicates is given below each
plot. Boxplots show the median (centre line) and the 25th and 75th
percentiles (box bounds); whiskers extend to the most extreme value
within 1.5 × the interquartile range of the nearest box bound, and all
plants are plotted as points. The statistics above the plots are
estimated marginal mean (EMM) contrasts of the BX effect within each
genotype (ΔEMM = BXplus − BXminus), computed from a linear mixed-effects
model with experiment as a random intercept to account for variation
between experiments. Each ΔEMM is a two-sided t-test on the estimated
marginal means (emmeans). ΔEMM is therefore expressed in the units of
the measurement, not in z-score units. (b) Summary of the same linear
mixed-effects model in (a). For each genotype, the BX effect is reported
as the EMM contrast Δ(BXplus − BXminus) with unadjusted two-sided P
values. The table also reports interaction contrasts relative to WT (ΔΔ
vs WT = Δ\_mutant − Δ\_WT), which test whether the BX response of each
mutant differs from that of WT; these two comparisons form one family
per measure and are two-sided t-tests with Holm adjusted P values. Data
are from Dataset 4.

\begin{longtable}[]{@{}llrrrrr@{}}
\caption{Estimated marginal means for mmf1 fresh weight by line and BX
condition (experiment included as a random intercept)}\tabularnewline
\toprule\noalign{}
line & BX & emmean & SE & df & lower.CL & upper.CL \\
\midrule\noalign{}
\endfirsthead
\toprule\noalign{}
line & BX & emmean & SE & df & lower.CL & upper.CL \\
\midrule\noalign{}
\endhead
\bottomrule\noalign{}
\endlastfoot
WT & minus & 0.265 & 0.088 & 5.081863 & 0.040 & 0.490 \\
mmf1-1 & minus & 0.208 & 0.088 & 5.233357 & -0.017 & 0.432 \\
mmf1-2 & minus & 0.244 & 0.088 & 5.108075 & 0.020 & 0.469 \\
WT & plus & 0.318 & 0.088 & 5.076884 & 0.093 & 0.542 \\
mmf1-1 & plus & 0.216 & 0.088 & 5.234357 & -0.008 & 0.440 \\
mmf1-2 & plus & 0.239 & 0.088 & 5.106284 & 0.015 & 0.464 \\
\end{longtable}

Extended Data Figure 5b (fresh weight) \textbar{} Growth feedback in a
BXplus microbiome is consistently lost in mmf1 mutants. (a and b)
Repeated growth feedbacks (Experiment III - V, VIII, IX and X) of Col-0
wild type (WT) and two mmf1 T-DNA insertion mutants (At3g44630) grown on
the BXminus and BXplus soil microbiomes. Feedback experiments were
performed on greenhouse-conditioned soil batches BS04 (Experiment III
and IV), BS09 (V and X), BS10 (IX) and BS11 (VIII). (a)
Z-score--transformed fresh weight and rosette area of WT and mmf1
mutants across six independent experiments. WT was present in all six
experiments, mmf1-1 in three and mmf1-2 in five. Values were
z-score--transformed within experiment to account for between-experiment
variation. Each point is an individual plant grown in its own pot,
i.e.~a biological replicate; number of replicates is given below each
plot. Boxplots show the median (centre line) and the 25th and 75th
percentiles (box bounds); whiskers extend to the most extreme value
within 1.5 × the interquartile range of the nearest box bound, and all
plants are plotted as points. The statistics above the plots are
estimated marginal mean (EMM) contrasts of the BX effect within each
genotype (ΔEMM = BXplus − BXminus), computed from a linear mixed-effects
model with experiment as a random intercept to account for variation
between experiments. Each ΔEMM is a two-sided t-test on the estimated
marginal means (emmeans). ΔEMM is therefore expressed in the units of
the measurement, not in z-score units. (b) Summary of the same linear
mixed-effects model in (a). For each genotype, the BX effect is reported
as the EMM contrast Δ(BXplus − BXminus) with unadjusted two-sided P
values. The table also reports interaction contrasts relative to WT (ΔΔ
vs WT = Δ\_mutant − Δ\_WT), which test whether the BX response of each
mutant differs from that of WT; these two comparisons form one family
per measure and are two-sided t-tests with Holm adjusted P values. Data
are from Dataset 4.

\begin{verbatim}
    line    BX emmean    SE       df lower.CL upper.CL
1     WT minus  0.265 0.088 5.081863    0.040    0.490
2 mmf1-1 minus  0.208 0.088 5.233357   -0.017    0.432
3 mmf1-2 minus  0.244 0.088 5.108075    0.020    0.469
4     WT  plus  0.318 0.088 5.076884    0.093    0.542
5 mmf1-1  plus  0.216 0.088 5.234357   -0.008    0.440
6 mmf1-2  plus  0.239 0.088 5.106284    0.015    0.464
\end{verbatim}

\begin{longtable}[]{@{}
  >{\raggedright\arraybackslash}p{(\linewidth - 14\tabcolsep) * \real{0.2766}}
  >{\raggedleft\arraybackslash}p{(\linewidth - 14\tabcolsep) * \real{0.1702}}
  >{\raggedleft\arraybackslash}p{(\linewidth - 14\tabcolsep) * \real{0.0638}}
  >{\raggedleft\arraybackslash}p{(\linewidth - 14\tabcolsep) * \real{0.0957}}
  >{\raggedleft\arraybackslash}p{(\linewidth - 14\tabcolsep) * \real{0.0957}}
  >{\raggedleft\arraybackslash}p{(\linewidth - 14\tabcolsep) * \real{0.0957}}
  >{\raggedleft\arraybackslash}p{(\linewidth - 14\tabcolsep) * \real{0.1064}}
  >{\raggedleft\arraybackslash}p{(\linewidth - 14\tabcolsep) * \real{0.0957}}@{}}
\caption{Fresh weight contrasts relative to WT in BXminus soil for mmf1
lines (experiment included as a random intercept;
Holm-adjusted)}\tabularnewline
\toprule\noalign{}
\begin{minipage}[b]{\linewidth}\raggedright
contrast
\end{minipage} & \begin{minipage}[b]{\linewidth}\raggedleft
diff\_vs\_WTminus
\end{minipage} & \begin{minipage}[b]{\linewidth}\raggedleft
SE
\end{minipage} & \begin{minipage}[b]{\linewidth}\raggedleft
df
\end{minipage} & \begin{minipage}[b]{\linewidth}\raggedleft
lower.CL
\end{minipage} & \begin{minipage}[b]{\linewidth}\raggedleft
upper.CL
\end{minipage} & \begin{minipage}[b]{\linewidth}\raggedleft
t
\end{minipage} & \begin{minipage}[b]{\linewidth}\raggedleft
p
\end{minipage} \\
\midrule\noalign{}
\endfirsthead
\toprule\noalign{}
\begin{minipage}[b]{\linewidth}\raggedright
contrast
\end{minipage} & \begin{minipage}[b]{\linewidth}\raggedleft
diff\_vs\_WTminus
\end{minipage} & \begin{minipage}[b]{\linewidth}\raggedleft
SE
\end{minipage} & \begin{minipage}[b]{\linewidth}\raggedleft
df
\end{minipage} & \begin{minipage}[b]{\linewidth}\raggedleft
lower.CL
\end{minipage} & \begin{minipage}[b]{\linewidth}\raggedleft
upper.CL
\end{minipage} & \begin{minipage}[b]{\linewidth}\raggedleft
t
\end{minipage} & \begin{minipage}[b]{\linewidth}\raggedleft
p
\end{minipage} \\
\midrule\noalign{}
\endhead
\bottomrule\noalign{}
\endlastfoot
(mmf1-1 minus) - WT minus & -0.057 & 0.017 & 374.1883 & -0.100 & -0.014
& -3.456509 & 0.002440 \\
(mmf1-2 minus) - WT minus & -0.021 & 0.014 & 374.0653 & -0.056 & 0.015 &
-1.517848 & 0.130000 \\
WT plus - WT minus & 0.053 & 0.013 & 374.0009 & 0.020 & 0.085 & 4.178160
& 0.000183 \\
(mmf1-1 plus) - WT minus & -0.049 & 0.017 & 374.1983 & -0.092 & -0.006 &
-2.951106 & 0.010100 \\
(mmf1-2 plus) - WT minus & -0.026 & 0.014 & 374.0681 & -0.061 & 0.009 &
-1.895243 & 0.118000 \\
\end{longtable}

Extended Data Figure 5b (fresh weight) \textbar{} Growth feedback in a
BXplus microbiome is consistently lost in mmf1 mutants. (a and b)
Repeated growth feedbacks (Experiment III - V, VIII, IX and X) of Col-0
wild type (WT) and two mmf1 T-DNA insertion mutants (At3g44630) grown on
the BXminus and BXplus soil microbiomes. Feedback experiments were
performed on greenhouse-conditioned soil batches BS04 (Experiment III
and IV), BS09 (V and X), BS10 (IX) and BS11 (VIII). (a)
Z-score--transformed fresh weight and rosette area of WT and mmf1
mutants across six independent experiments. WT was present in all six
experiments, mmf1-1 in three and mmf1-2 in five. Values were
z-score--transformed within experiment to account for between-experiment
variation. Each point is an individual plant grown in its own pot,
i.e.~a biological replicate; number of replicates is given below each
plot. Boxplots show the median (centre line) and the 25th and 75th
percentiles (box bounds); whiskers extend to the most extreme value
within 1.5 × the interquartile range of the nearest box bound, and all
plants are plotted as points. The statistics above the plots are
estimated marginal mean (EMM) contrasts of the BX effect within each
genotype (ΔEMM = BXplus − BXminus), computed from a linear mixed-effects
model with experiment as a random intercept to account for variation
between experiments. Each ΔEMM is a two-sided t-test on the estimated
marginal means (emmeans). ΔEMM is therefore expressed in the units of
the measurement, not in z-score units. (b) Summary of the same linear
mixed-effects model in (a). For each genotype, the BX effect is reported
as the EMM contrast Δ(BXplus − BXminus) with unadjusted two-sided P
values. The table also reports interaction contrasts relative to WT (ΔΔ
vs WT = Δ\_mutant − Δ\_WT), which test whether the BX response of each
mutant differs from that of WT; these two comparisons form one family
per measure and are two-sided t-tests with Holm adjusted P values. Data
are from Dataset 4.

\subsubsection{\texorpdfstring{\textbf{Plot combined mmf1 fresh
weight}}{Plot combined mmf1 fresh weight}}\label{plot-combined-mmf1-fresh-weight}

\begin{figure}

\centering{

\pandocbounded{\includegraphics[keepaspectratio]{04_Mutant_screening_files/figure-pdf/fig-mmf1-fw-combined-1.pdf}}

}

\caption{\label{fig-mmf1-fw-combined}Extended Data Figure 5a (fresh
weight) \textbar{} Growth feedback in a BXplus microbiome is
consistently lost in mmf1 mutants. (a and b) Repeated growth feedbacks
(Experiment III - V, VIII, IX and X) of Col-0 wild type (WT) and two
mmf1 T-DNA insertion mutants (At3g44630) grown on the BXminus and BXplus
soil microbiomes. Feedback experiments were performed on
greenhouse-conditioned soil batches BS04 (Experiment III and IV), BS09
(V and X), BS10 (IX) and BS11 (VIII). (a) Z-score--transformed fresh
weight and rosette area of WT and mmf1 mutants across six independent
experiments. WT was present in all six experiments, mmf1-1 in three and
mmf1-2 in five. Values were z-score--transformed within experiment to
account for between-experiment variation. Each point is an individual
plant grown in its own pot, i.e.~a biological replicate; number of
replicates is given below each plot. Boxplots show the median (centre
line) and the 25th and 75th percentiles (box bounds); whiskers extend to
the most extreme value within 1.5 × the interquartile range of the
nearest box bound, and all plants are plotted as points. The statistics
above the plots are estimated marginal mean (EMM) contrasts of the BX
effect within each genotype (ΔEMM = BXplus − BXminus), computed from a
linear mixed-effects model with experiment as a random intercept to
account for variation between experiments. Each ΔEMM is a two-sided
t-test on the estimated marginal means (emmeans). ΔEMM is therefore
expressed in the units of the measurement, not in z-score units. (b)
Summary of the same linear mixed-effects model in (a). For each
genotype, the BX effect is reported as the EMM contrast Δ(BXplus −
BXminus) with unadjusted two-sided P values. The table also reports
interaction contrasts relative to WT (ΔΔ vs WT = Δ\_mutant − Δ\_WT),
which test whether the BX response of each mutant differs from that of
WT; these two comparisons form one family per measure and are two-sided
t-tests with Holm adjusted P values. Data are from Dataset 4.}

\end{figure}%

\subsubsection{LMM and EMM:rosette area}\label{lmm-and-emmrosette-area}

\begin{verbatim}
Linear mixed model fit by REML. t-tests use Satterthwaite's method [
lmerModLmerTest]
Formula: formula_lmm
   Data: data

REML criterion at convergence: 10489.7

Scaled residuals: 
    Min      1Q  Median      3Q     Max 
-2.7775 -0.5483 -0.0011  0.4903  3.9084 

Random effects:
 Groups     Name        Variance Std.Dev.
 experiment (Intercept) 1.14e+11 337586  
 Residual               5.39e+10 232174  
Number of obs: 385, groups:  experiment, 6

Fixed effects:
                    Estimate Std. Error         df t value Pr(>|t|)    
(Intercept)        778742.15  140186.33       5.27   5.555  0.00220 ** 
BXplus             150434.70   35844.22     374.01   4.197 3.38e-05 ***
linemmf1-1        -290807.29   47084.30     374.60  -6.176 1.71e-09 ***
linemmf1-2         -94129.22   38461.19     374.21  -2.447  0.01485 *  
BXplus:linemmf1-1  -93502.32   63638.23     374.01  -1.469  0.14260    
BXplus:linemmf1-2 -162681.39   53272.47     374.01  -3.054  0.00242 ** 
---
Signif. codes:  0 '***' 0.001 '**' 0.01 '*' 0.05 '.' 0.1 ' ' 1

Correlation of Fixed Effects:
            (Intr) BXplus lnm1-1 lnm1-2 BX:1-1
BXplus      -0.131                            
linemmf1-1  -0.098  0.386                     
linemmf1-2  -0.122  0.477  0.330              
BXpls:ln1-1  0.074 -0.563 -0.678 -0.271       
BXpls:ln1-2  0.088 -0.673 -0.261 -0.702  0.379
\end{verbatim}

\begin{longtable}[]{@{}
  >{\raggedright\arraybackslash}p{(\linewidth - 16\tabcolsep) * \real{0.1354}}
  >{\raggedright\arraybackslash}p{(\linewidth - 16\tabcolsep) * \real{0.0729}}
  >{\raggedleft\arraybackslash}p{(\linewidth - 16\tabcolsep) * \real{0.1875}}
  >{\raggedleft\arraybackslash}p{(\linewidth - 16\tabcolsep) * \real{0.0938}}
  >{\raggedleft\arraybackslash}p{(\linewidth - 16\tabcolsep) * \real{0.0938}}
  >{\raggedleft\arraybackslash}p{(\linewidth - 16\tabcolsep) * \real{0.1042}}
  >{\raggedleft\arraybackslash}p{(\linewidth - 16\tabcolsep) * \real{0.1042}}
  >{\raggedleft\arraybackslash}p{(\linewidth - 16\tabcolsep) * \real{0.1146}}
  >{\raggedleft\arraybackslash}p{(\linewidth - 16\tabcolsep) * \real{0.0938}}@{}}
\caption{Within-line ΔEMM = BXplus - BXminus for mmf1 rosette
area}\tabularnewline
\toprule\noalign{}
\begin{minipage}[b]{\linewidth}\raggedright
contrast
\end{minipage} & \begin{minipage}[b]{\linewidth}\raggedright
line
\end{minipage} & \begin{minipage}[b]{\linewidth}\raggedleft
effect\_plus\_minus
\end{minipage} & \begin{minipage}[b]{\linewidth}\raggedleft
SE
\end{minipage} & \begin{minipage}[b]{\linewidth}\raggedleft
df
\end{minipage} & \begin{minipage}[b]{\linewidth}\raggedleft
lower.CL
\end{minipage} & \begin{minipage}[b]{\linewidth}\raggedleft
upper.CL
\end{minipage} & \begin{minipage}[b]{\linewidth}\raggedleft
t
\end{minipage} & \begin{minipage}[b]{\linewidth}\raggedleft
p
\end{minipage} \\
\midrule\noalign{}
\endfirsthead
\toprule\noalign{}
\begin{minipage}[b]{\linewidth}\raggedright
contrast
\end{minipage} & \begin{minipage}[b]{\linewidth}\raggedright
line
\end{minipage} & \begin{minipage}[b]{\linewidth}\raggedleft
effect\_plus\_minus
\end{minipage} & \begin{minipage}[b]{\linewidth}\raggedleft
SE
\end{minipage} & \begin{minipage}[b]{\linewidth}\raggedleft
df
\end{minipage} & \begin{minipage}[b]{\linewidth}\raggedleft
lower.CL
\end{minipage} & \begin{minipage}[b]{\linewidth}\raggedleft
upper.CL
\end{minipage} & \begin{minipage}[b]{\linewidth}\raggedleft
t
\end{minipage} & \begin{minipage}[b]{\linewidth}\raggedleft
p
\end{minipage} \\
\midrule\noalign{}
\endhead
\bottomrule\noalign{}
\endlastfoot
plus - minus & WT & 150434.70 & 35844.27 & 374.0029 & 79953.14 &
220916.27 & 4.1968966 & 3.38e-05 \\
plus - minus & mmf1-1 & 56932.38 & 52587.13 & 374.0022 & -46471.12 &
160335.89 & 1.0826295 & 2.80e-01 \\
plus - minus & mmf1-2 & -12246.68 & 39412.85 & 374.0083 & -89745.24 &
65251.88 & -0.3107281 & 7.56e-01 \\
\end{longtable}

Extended Data Figure 5b (rosette area) \textbar{} Growth feedback in a
BXplus microbiome is consistently lost in mmf1 mutants. (a and b)
Repeated growth feedbacks (Experiment III - V, VIII, IX and X) of Col-0
wild type (WT) and two mmf1 T-DNA insertion mutants (At3g44630) grown on
the BXminus and BXplus soil microbiomes. Feedback experiments were
performed on greenhouse-conditioned soil batches BS04 (Experiment III
and IV), BS09 (V and X), BS10 (IX) and BS11 (VIII). (a)
Z-score--transformed fresh weight and rosette area of WT and mmf1
mutants across six independent experiments. WT was present in all six
experiments, mmf1-1 in three and mmf1-2 in five. Values were
z-score--transformed within experiment to account for between-experiment
variation. Each point is an individual plant grown in its own pot,
i.e.~a biological replicate; number of replicates is given below each
plot. Boxplots show the median (centre line) and the 25th and 75th
percentiles (box bounds); whiskers extend to the most extreme value
within 1.5 × the interquartile range of the nearest box bound, and all
plants are plotted as points. The statistics above the plots are
estimated marginal mean (EMM) contrasts of the BX effect within each
genotype (ΔEMM = BXplus − BXminus), computed from a linear mixed-effects
model with experiment as a random intercept to account for variation
between experiments. Each ΔEMM is a two-sided t-test on the estimated
marginal means (emmeans). ΔEMM is therefore expressed in the units of
the measurement, not in z-score units. (b) Summary of the same linear
mixed-effects model in (a). For each genotype, the BX effect is reported
as the EMM contrast Δ(BXplus − BXminus) with unadjusted two-sided P
values. The table also reports interaction contrasts relative to WT (ΔΔ
vs WT = Δ\_mutant − Δ\_WT), which test whether the BX response of each
mutant differs from that of WT; these two comparisons form one family
per measure and are two-sided t-tests with Holm adjusted P values. Data
are from Dataset 4.

\begin{longtable}[]{@{}
  >{\raggedright\arraybackslash}p{(\linewidth - 16\tabcolsep) * \real{0.1515}}
  >{\raggedright\arraybackslash}p{(\linewidth - 16\tabcolsep) * \real{0.1717}}
  >{\raggedleft\arraybackslash}p{(\linewidth - 16\tabcolsep) * \real{0.1111}}
  >{\raggedleft\arraybackslash}p{(\linewidth - 16\tabcolsep) * \real{0.0909}}
  >{\raggedleft\arraybackslash}p{(\linewidth - 16\tabcolsep) * \real{0.0909}}
  >{\raggedleft\arraybackslash}p{(\linewidth - 16\tabcolsep) * \real{0.1010}}
  >{\raggedleft\arraybackslash}p{(\linewidth - 16\tabcolsep) * \real{0.1010}}
  >{\raggedleft\arraybackslash}p{(\linewidth - 16\tabcolsep) * \real{0.1010}}
  >{\raggedleft\arraybackslash}p{(\linewidth - 16\tabcolsep) * \real{0.0808}}@{}}
\caption{Difference-in-differences relative to WT for mmf1 rosette
area}\tabularnewline
\toprule\noalign{}
\begin{minipage}[b]{\linewidth}\raggedright
BX\_revpairwise
\end{minipage} & \begin{minipage}[b]{\linewidth}\raggedright
line\_trt.vs.ctrl
\end{minipage} & \begin{minipage}[b]{\linewidth}\raggedleft
diff\_vs\_WT
\end{minipage} & \begin{minipage}[b]{\linewidth}\raggedleft
SE
\end{minipage} & \begin{minipage}[b]{\linewidth}\raggedleft
df
\end{minipage} & \begin{minipage}[b]{\linewidth}\raggedleft
lower.CL
\end{minipage} & \begin{minipage}[b]{\linewidth}\raggedleft
upper.CL
\end{minipage} & \begin{minipage}[b]{\linewidth}\raggedleft
t
\end{minipage} & \begin{minipage}[b]{\linewidth}\raggedleft
p\_adj
\end{minipage} \\
\midrule\noalign{}
\endfirsthead
\toprule\noalign{}
\begin{minipage}[b]{\linewidth}\raggedright
BX\_revpairwise
\end{minipage} & \begin{minipage}[b]{\linewidth}\raggedright
line\_trt.vs.ctrl
\end{minipage} & \begin{minipage}[b]{\linewidth}\raggedleft
diff\_vs\_WT
\end{minipage} & \begin{minipage}[b]{\linewidth}\raggedleft
SE
\end{minipage} & \begin{minipage}[b]{\linewidth}\raggedleft
df
\end{minipage} & \begin{minipage}[b]{\linewidth}\raggedleft
lower.CL
\end{minipage} & \begin{minipage}[b]{\linewidth}\raggedleft
upper.CL
\end{minipage} & \begin{minipage}[b]{\linewidth}\raggedleft
t
\end{minipage} & \begin{minipage}[b]{\linewidth}\raggedleft
p\_adj
\end{minipage} \\
\midrule\noalign{}
\endhead
\bottomrule\noalign{}
\endlastfoot
plus - minus & (mmf1-1) - WT & -93502.32 & 63638.31 & 374.0022 &
-236718.0 & 49713.34 & -1.469277 & 0.14300 \\
plus - minus & (mmf1-2) - WT & -162681.39 & 53272.63 & 374.0053 &
-282569.5 & -42793.31 & -3.053752 & 0.00484 \\
\end{longtable}

Extended Data Figure 5b (rosette area) \textbar{} Growth feedback in a
BXplus microbiome is consistently lost in mmf1 mutants. (a and b)
Repeated growth feedbacks (Experiment III - V, VIII, IX and X) of Col-0
wild type (WT) and two mmf1 T-DNA insertion mutants (At3g44630) grown on
the BXminus and BXplus soil microbiomes. Feedback experiments were
performed on greenhouse-conditioned soil batches BS04 (Experiment III
and IV), BS09 (V and X), BS10 (IX) and BS11 (VIII). (a)
Z-score--transformed fresh weight and rosette area of WT and mmf1
mutants across six independent experiments. WT was present in all six
experiments, mmf1-1 in three and mmf1-2 in five. Values were
z-score--transformed within experiment to account for between-experiment
variation. Each point is an individual plant grown in its own pot,
i.e.~a biological replicate; number of replicates is given below each
plot. Boxplots show the median (centre line) and the 25th and 75th
percentiles (box bounds); whiskers extend to the most extreme value
within 1.5 × the interquartile range of the nearest box bound, and all
plants are plotted as points. The statistics above the plots are
estimated marginal mean (EMM) contrasts of the BX effect within each
genotype (ΔEMM = BXplus − BXminus), computed from a linear mixed-effects
model with experiment as a random intercept to account for variation
between experiments. Each ΔEMM is a two-sided t-test on the estimated
marginal means (emmeans). ΔEMM is therefore expressed in the units of
the measurement, not in z-score units. (b) Summary of the same linear
mixed-effects model in (a). For each genotype, the BX effect is reported
as the EMM contrast Δ(BXplus − BXminus) with unadjusted two-sided P
values. The table also reports interaction contrasts relative to WT (ΔΔ
vs WT = Δ\_mutant − Δ\_WT), which test whether the BX response of each
mutant differs from that of WT; these two comparisons form one family
per measure and are two-sided t-tests with Holm adjusted P values. Data
are from Dataset 4.

\begin{longtable}[]{@{}
  >{\raggedright\arraybackslash}p{(\linewidth - 16\tabcolsep) * \real{0.1573}}
  >{\raggedright\arraybackslash}p{(\linewidth - 16\tabcolsep) * \real{0.0674}}
  >{\raggedleft\arraybackslash}p{(\linewidth - 16\tabcolsep) * \real{0.1236}}
  >{\raggedleft\arraybackslash}p{(\linewidth - 16\tabcolsep) * \real{0.1011}}
  >{\raggedleft\arraybackslash}p{(\linewidth - 16\tabcolsep) * \real{0.1011}}
  >{\raggedleft\arraybackslash}p{(\linewidth - 16\tabcolsep) * \real{0.1124}}
  >{\raggedleft\arraybackslash}p{(\linewidth - 16\tabcolsep) * \real{0.1348}}
  >{\raggedleft\arraybackslash}p{(\linewidth - 16\tabcolsep) * \real{0.1124}}
  >{\raggedleft\arraybackslash}p{(\linewidth - 16\tabcolsep) * \real{0.0899}}@{}}
\caption{Mutant vs WT contrasts within each BX condition for mmf1
rosette area (Holm-adjusted)}\tabularnewline
\toprule\noalign{}
\begin{minipage}[b]{\linewidth}\raggedright
contrast
\end{minipage} & \begin{minipage}[b]{\linewidth}\raggedright
BX
\end{minipage} & \begin{minipage}[b]{\linewidth}\raggedleft
estimate
\end{minipage} & \begin{minipage}[b]{\linewidth}\raggedleft
SE
\end{minipage} & \begin{minipage}[b]{\linewidth}\raggedleft
df
\end{minipage} & \begin{minipage}[b]{\linewidth}\raggedleft
lower.CL
\end{minipage} & \begin{minipage}[b]{\linewidth}\raggedleft
upper.CL
\end{minipage} & \begin{minipage}[b]{\linewidth}\raggedleft
t.ratio
\end{minipage} & \begin{minipage}[b]{\linewidth}\raggedleft
p.value
\end{minipage} \\
\midrule\noalign{}
\endfirsthead
\toprule\noalign{}
\begin{minipage}[b]{\linewidth}\raggedright
contrast
\end{minipage} & \begin{minipage}[b]{\linewidth}\raggedright
BX
\end{minipage} & \begin{minipage}[b]{\linewidth}\raggedleft
estimate
\end{minipage} & \begin{minipage}[b]{\linewidth}\raggedleft
SE
\end{minipage} & \begin{minipage}[b]{\linewidth}\raggedleft
df
\end{minipage} & \begin{minipage}[b]{\linewidth}\raggedleft
lower.CL
\end{minipage} & \begin{minipage}[b]{\linewidth}\raggedleft
upper.CL
\end{minipage} & \begin{minipage}[b]{\linewidth}\raggedleft
t.ratio
\end{minipage} & \begin{minipage}[b]{\linewidth}\raggedleft
p.value
\end{minipage} \\
\midrule\noalign{}
\endhead
\bottomrule\noalign{}
\endlastfoot
(mmf1-1) - WT & minus & -290807.29 & 47099.84 & 374.5925 & -396803.1 &
-184811.512 & -6.174273 & 0.0000 \\
(mmf1-2) - WT & minus & -94129.22 & 38465.56 & 374.2082 & -180694.4 &
-7564.089 & -2.447104 & 0.0149 \\
(mmf1-1) - WT & plus & -384309.61 & 46987.51 & 374.6677 & -490052.5 &
-278566.713 & -8.178974 & 0.0000 \\
(mmf1-2) - WT & plus & -256810.61 & 37934.27 & 374.2153 & -342180.1 &
-171441.121 & -6.769883 & 0.0000 \\
\end{longtable}

Extended Data Figure 5b (rosette area) \textbar{} Growth feedback in a
BXplus microbiome is consistently lost in mmf1 mutants. (a and b)
Repeated growth feedbacks (Experiment III - V, VIII, IX and X) of Col-0
wild type (WT) and two mmf1 T-DNA insertion mutants (At3g44630) grown on
the BXminus and BXplus soil microbiomes. Feedback experiments were
performed on greenhouse-conditioned soil batches BS04 (Experiment III
and IV), BS09 (V and X), BS10 (IX) and BS11 (VIII). (a)
Z-score--transformed fresh weight and rosette area of WT and mmf1
mutants across six independent experiments. WT was present in all six
experiments, mmf1-1 in three and mmf1-2 in five. Values were
z-score--transformed within experiment to account for between-experiment
variation. Each point is an individual plant grown in its own pot,
i.e.~a biological replicate; number of replicates is given below each
plot. Boxplots show the median (centre line) and the 25th and 75th
percentiles (box bounds); whiskers extend to the most extreme value
within 1.5 × the interquartile range of the nearest box bound, and all
plants are plotted as points. The statistics above the plots are
estimated marginal mean (EMM) contrasts of the BX effect within each
genotype (ΔEMM = BXplus − BXminus), computed from a linear mixed-effects
model with experiment as a random intercept to account for variation
between experiments. Each ΔEMM is a two-sided t-test on the estimated
marginal means (emmeans). ΔEMM is therefore expressed in the units of
the measurement, not in z-score units. (b) Summary of the same linear
mixed-effects model in (a). For each genotype, the BX effect is reported
as the EMM contrast Δ(BXplus − BXminus) with unadjusted two-sided P
values. The table also reports interaction contrasts relative to WT (ΔΔ
vs WT = Δ\_mutant − Δ\_WT), which test whether the BX response of each
mutant differs from that of WT; these two comparisons form one family
per measure and are two-sided t-tests with Holm adjusted P values. Data
are from Dataset 4.

\begin{longtable}[]{@{}lrrrrr@{}}
\caption{Estimated marginal means for mmf1 rosette area by line,
averaged across BX conditions (experiment included as a random
intercept)}\tabularnewline
\toprule\noalign{}
line & emmean & SE & df & lower.CL & upper.CL \\
\midrule\noalign{}
\endfirsthead
\toprule\noalign{}
line & emmean & SE & df & lower.CL & upper.CL \\
\midrule\noalign{}
\endhead
\bottomrule\noalign{}
\endlastfoot
WT & 853959.5 & 138981.2 & 5.087205 & 498531.1 & 1209387.9 \\
mmf1-1 & 516401.1 & 141012.6 & 5.389086 & 161648.5 & 871153.6 \\
mmf1-2 & 678489.6 & 139350.9 & 5.141190 & 323215.5 & 1033763.7 \\
\end{longtable}

Extended Data Figure 5b (rosette area) \textbar{} Growth feedback in a
BXplus microbiome is consistently lost in mmf1 mutants. (a and b)
Repeated growth feedbacks (Experiment III - V, VIII, IX and X) of Col-0
wild type (WT) and two mmf1 T-DNA insertion mutants (At3g44630) grown on
the BXminus and BXplus soil microbiomes. Feedback experiments were
performed on greenhouse-conditioned soil batches BS04 (Experiment III
and IV), BS09 (V and X), BS10 (IX) and BS11 (VIII). (a)
Z-score--transformed fresh weight and rosette area of WT and mmf1
mutants across six independent experiments. WT was present in all six
experiments, mmf1-1 in three and mmf1-2 in five. Values were
z-score--transformed within experiment to account for between-experiment
variation. Each point is an individual plant grown in its own pot,
i.e.~a biological replicate; number of replicates is given below each
plot. Boxplots show the median (centre line) and the 25th and 75th
percentiles (box bounds); whiskers extend to the most extreme value
within 1.5 × the interquartile range of the nearest box bound, and all
plants are plotted as points. The statistics above the plots are
estimated marginal mean (EMM) contrasts of the BX effect within each
genotype (ΔEMM = BXplus − BXminus), computed from a linear mixed-effects
model with experiment as a random intercept to account for variation
between experiments. Each ΔEMM is a two-sided t-test on the estimated
marginal means (emmeans). ΔEMM is therefore expressed in the units of
the measurement, not in z-score units. (b) Summary of the same linear
mixed-effects model in (a). For each genotype, the BX effect is reported
as the EMM contrast Δ(BXplus − BXminus) with unadjusted two-sided P
values. The table also reports interaction contrasts relative to WT (ΔΔ
vs WT = Δ\_mutant − Δ\_WT), which test whether the BX response of each
mutant differs from that of WT; these two comparisons form one family
per measure and are two-sided t-tests with Holm adjusted P values. Data
are from Dataset 4.

\begin{longtable}[]{@{}
  >{\raggedright\arraybackslash}p{(\linewidth - 14\tabcolsep) * \real{0.1867}}
  >{\raggedleft\arraybackslash}p{(\linewidth - 14\tabcolsep) * \real{0.1333}}
  >{\raggedleft\arraybackslash}p{(\linewidth - 14\tabcolsep) * \real{0.1200}}
  >{\raggedleft\arraybackslash}p{(\linewidth - 14\tabcolsep) * \real{0.1200}}
  >{\raggedleft\arraybackslash}p{(\linewidth - 14\tabcolsep) * \real{0.1333}}
  >{\raggedleft\arraybackslash}p{(\linewidth - 14\tabcolsep) * \real{0.1333}}
  >{\raggedleft\arraybackslash}p{(\linewidth - 14\tabcolsep) * \real{0.1333}}
  >{\raggedleft\arraybackslash}p{(\linewidth - 14\tabcolsep) * \real{0.0400}}@{}}
\caption{Mutant vs WT contrasts for mmf1 rosette area, averaged across
BX conditions and corrected for experiment
(Holm-adjusted)}\tabularnewline
\toprule\noalign{}
\begin{minipage}[b]{\linewidth}\raggedright
contrast
\end{minipage} & \begin{minipage}[b]{\linewidth}\raggedleft
diff\_EMM
\end{minipage} & \begin{minipage}[b]{\linewidth}\raggedleft
SE
\end{minipage} & \begin{minipage}[b]{\linewidth}\raggedleft
df
\end{minipage} & \begin{minipage}[b]{\linewidth}\raggedleft
lower.CL
\end{minipage} & \begin{minipage}[b]{\linewidth}\raggedleft
upper.CL
\end{minipage} & \begin{minipage}[b]{\linewidth}\raggedleft
t
\end{minipage} & \begin{minipage}[b]{\linewidth}\raggedleft
p
\end{minipage} \\
\midrule\noalign{}
\endfirsthead
\toprule\noalign{}
\begin{minipage}[b]{\linewidth}\raggedright
contrast
\end{minipage} & \begin{minipage}[b]{\linewidth}\raggedleft
diff\_EMM
\end{minipage} & \begin{minipage}[b]{\linewidth}\raggedleft
SE
\end{minipage} & \begin{minipage}[b]{\linewidth}\raggedleft
df
\end{minipage} & \begin{minipage}[b]{\linewidth}\raggedleft
lower.CL
\end{minipage} & \begin{minipage}[b]{\linewidth}\raggedleft
upper.CL
\end{minipage} & \begin{minipage}[b]{\linewidth}\raggedleft
t
\end{minipage} & \begin{minipage}[b]{\linewidth}\raggedleft
p
\end{minipage} \\
\midrule\noalign{}
\endhead
\bottomrule\noalign{}
\endlastfoot
(mmf1-1) - WT & -337558.5 & 34650.42 & 375.1285 & -415537.0 & -259579.9
& -9.741827 & 0 \\
(mmf1-2) - WT & -175469.9 & 27382.68 & 374.4032 & -237093.4 & -113846.5
& -6.408063 & 0 \\
\end{longtable}

Extended Data Figure 5b (rosette area) \textbar{} Growth feedback in a
BXplus microbiome is consistently lost in mmf1 mutants. (a and b)
Repeated growth feedbacks (Experiment III - V, VIII, IX and X) of Col-0
wild type (WT) and two mmf1 T-DNA insertion mutants (At3g44630) grown on
the BXminus and BXplus soil microbiomes. Feedback experiments were
performed on greenhouse-conditioned soil batches BS04 (Experiment III
and IV), BS09 (V and X), BS10 (IX) and BS11 (VIII). (a)
Z-score--transformed fresh weight and rosette area of WT and mmf1
mutants across six independent experiments. WT was present in all six
experiments, mmf1-1 in three and mmf1-2 in five. Values were
z-score--transformed within experiment to account for between-experiment
variation. Each point is an individual plant grown in its own pot,
i.e.~a biological replicate; number of replicates is given below each
plot. Boxplots show the median (centre line) and the 25th and 75th
percentiles (box bounds); whiskers extend to the most extreme value
within 1.5 × the interquartile range of the nearest box bound, and all
plants are plotted as points. The statistics above the plots are
estimated marginal mean (EMM) contrasts of the BX effect within each
genotype (ΔEMM = BXplus − BXminus), computed from a linear mixed-effects
model with experiment as a random intercept to account for variation
between experiments. Each ΔEMM is a two-sided t-test on the estimated
marginal means (emmeans). ΔEMM is therefore expressed in the units of
the measurement, not in z-score units. (b) Summary of the same linear
mixed-effects model in (a). For each genotype, the BX effect is reported
as the EMM contrast Δ(BXplus − BXminus) with unadjusted two-sided P
values. The table also reports interaction contrasts relative to WT (ΔΔ
vs WT = Δ\_mutant − Δ\_WT), which test whether the BX response of each
mutant differs from that of WT; these two comparisons form one family
per measure and are two-sided t-tests with Holm adjusted P values. Data
are from Dataset 4.

\begin{longtable}[]{@{}llrrrrr@{}}
\caption{Estimated marginal means for mmf1 rosette area by line and BX
condition (experiment included as a random intercept)}\tabularnewline
\toprule\noalign{}
line & BX & emmean & SE & df & lower.CL & upper.CL \\
\midrule\noalign{}
\endfirsthead
\toprule\noalign{}
line & BX & emmean & SE & df & lower.CL & upper.CL \\
\midrule\noalign{}
\endhead
\bottomrule\noalign{}
\endlastfoot
WT & minus & 778742.2 & 140186.4 & 5.265916 & 423786.8 & 1133697.5 \\
mmf1-1 & minus & 487934.9 & 143432.5 & 5.768409 & 133524.1 & 842345.6 \\
mmf1-2 & minus & 684612.9 & 140756.8 & 5.351760 & 329823.5 &
1039402.3 \\
WT & plus & 929176.9 & 140077.7 & 5.249612 & 574185.0 & 1284168.7 \\
mmf1-1 & plus & 544867.2 & 143453.6 & 5.771703 & 190455.6 & 899278.9 \\
mmf1-2 & plus & 672366.2 & 140718.0 & 5.345855 & 317566.0 & 1027166.5 \\
\end{longtable}

Extended Data Figure 5b (rosette area) \textbar{} Growth feedback in a
BXplus microbiome is consistently lost in mmf1 mutants. (a and b)
Repeated growth feedbacks (Experiment III - V, VIII, IX and X) of Col-0
wild type (WT) and two mmf1 T-DNA insertion mutants (At3g44630) grown on
the BXminus and BXplus soil microbiomes. Feedback experiments were
performed on greenhouse-conditioned soil batches BS04 (Experiment III
and IV), BS09 (V and X), BS10 (IX) and BS11 (VIII). (a)
Z-score--transformed fresh weight and rosette area of WT and mmf1
mutants across six independent experiments. WT was present in all six
experiments, mmf1-1 in three and mmf1-2 in five. Values were
z-score--transformed within experiment to account for between-experiment
variation. Each point is an individual plant grown in its own pot,
i.e.~a biological replicate; number of replicates is given below each
plot. Boxplots show the median (centre line) and the 25th and 75th
percentiles (box bounds); whiskers extend to the most extreme value
within 1.5 × the interquartile range of the nearest box bound, and all
plants are plotted as points. The statistics above the plots are
estimated marginal mean (EMM) contrasts of the BX effect within each
genotype (ΔEMM = BXplus − BXminus), computed from a linear mixed-effects
model with experiment as a random intercept to account for variation
between experiments. Each ΔEMM is a two-sided t-test on the estimated
marginal means (emmeans). ΔEMM is therefore expressed in the units of
the measurement, not in z-score units. (b) Summary of the same linear
mixed-effects model in (a). For each genotype, the BX effect is reported
as the EMM contrast Δ(BXplus − BXminus) with unadjusted two-sided P
values. The table also reports interaction contrasts relative to WT (ΔΔ
vs WT = Δ\_mutant − Δ\_WT), which test whether the BX response of each
mutant differs from that of WT; these two comparisons form one family
per measure and are two-sided t-tests with Holm adjusted P values. Data
are from Dataset 4.

\begin{verbatim}
    line    BX   emmean       SE       df lower.CL  upper.CL
1     WT minus 778742.2 140186.4 5.265916 423786.8 1133697.5
2 mmf1-1 minus 487934.9 143432.5 5.768409 133524.1  842345.6
3 mmf1-2 minus 684612.9 140756.8 5.351760 329823.5 1039402.3
4     WT  plus 929176.9 140077.7 5.249612 574185.0 1284168.7
5 mmf1-1  plus 544867.2 143453.6 5.771703 190455.6  899278.9
6 mmf1-2  plus 672366.2 140718.0 5.345855 317566.0 1027166.5
\end{verbatim}

\begin{longtable}[]{@{}
  >{\raggedright\arraybackslash}p{(\linewidth - 14\tabcolsep) * \real{0.2549}}
  >{\raggedleft\arraybackslash}p{(\linewidth - 14\tabcolsep) * \real{0.1569}}
  >{\raggedleft\arraybackslash}p{(\linewidth - 14\tabcolsep) * \real{0.0882}}
  >{\raggedleft\arraybackslash}p{(\linewidth - 14\tabcolsep) * \real{0.0882}}
  >{\raggedleft\arraybackslash}p{(\linewidth - 14\tabcolsep) * \real{0.1078}}
  >{\raggedleft\arraybackslash}p{(\linewidth - 14\tabcolsep) * \real{0.1176}}
  >{\raggedleft\arraybackslash}p{(\linewidth - 14\tabcolsep) * \real{0.0980}}
  >{\raggedleft\arraybackslash}p{(\linewidth - 14\tabcolsep) * \real{0.0882}}@{}}
\caption{Rosette area contrasts relative to WT in BXminus soil for mmf1
lines (experiment included as a random intercept;
Holm-adjusted)}\tabularnewline
\toprule\noalign{}
\begin{minipage}[b]{\linewidth}\raggedright
contrast
\end{minipage} & \begin{minipage}[b]{\linewidth}\raggedleft
diff\_vs\_WTminus
\end{minipage} & \begin{minipage}[b]{\linewidth}\raggedleft
SE
\end{minipage} & \begin{minipage}[b]{\linewidth}\raggedleft
df
\end{minipage} & \begin{minipage}[b]{\linewidth}\raggedleft
lower.CL
\end{minipage} & \begin{minipage}[b]{\linewidth}\raggedleft
upper.CL
\end{minipage} & \begin{minipage}[b]{\linewidth}\raggedleft
t
\end{minipage} & \begin{minipage}[b]{\linewidth}\raggedleft
p
\end{minipage} \\
\midrule\noalign{}
\endfirsthead
\toprule\noalign{}
\begin{minipage}[b]{\linewidth}\raggedright
contrast
\end{minipage} & \begin{minipage}[b]{\linewidth}\raggedleft
diff\_vs\_WTminus
\end{minipage} & \begin{minipage}[b]{\linewidth}\raggedleft
SE
\end{minipage} & \begin{minipage}[b]{\linewidth}\raggedleft
df
\end{minipage} & \begin{minipage}[b]{\linewidth}\raggedleft
lower.CL
\end{minipage} & \begin{minipage}[b]{\linewidth}\raggedleft
upper.CL
\end{minipage} & \begin{minipage}[b]{\linewidth}\raggedleft
t
\end{minipage} & \begin{minipage}[b]{\linewidth}\raggedleft
p
\end{minipage} \\
\midrule\noalign{}
\endhead
\bottomrule\noalign{}
\endlastfoot
(mmf1-1 minus) - WT minus & -290807.29 & 47099.84 & 374.5925 &
-412749.60 & -168864.978 & -6.174273 & 0.00e+00 \\
(mmf1-2 minus) - WT minus & -94129.22 & 38465.56 & 374.2082 & -193717.75
& 5459.312 & -2.447104 & 1.49e-02 \\
WT plus - WT minus & 150434.70 & 35844.27 & 374.0029 & 57632.51 &
243236.903 & 4.196897 & 1.02e-04 \\
(mmf1-1 plus) - WT minus & -233874.91 & 47166.40 & 374.6229 & -355989.50
& -111760.315 & -4.958506 & 4.30e-06 \\
(mmf1-2 plus) - WT minus & -106375.90 & 38333.83 & 374.2170 & -205623.36
& -7128.445 & -2.774988 & 1.16e-02 \\
\end{longtable}

Extended Data Figure 5b (rosette area) \textbar{} Growth feedback in a
BXplus microbiome is consistently lost in mmf1 mutants. (a and b)
Repeated growth feedbacks (Experiment III - V, VIII, IX and X) of Col-0
wild type (WT) and two mmf1 T-DNA insertion mutants (At3g44630) grown on
the BXminus and BXplus soil microbiomes. Feedback experiments were
performed on greenhouse-conditioned soil batches BS04 (Experiment III
and IV), BS09 (V and X), BS10 (IX) and BS11 (VIII). (a)
Z-score--transformed fresh weight and rosette area of WT and mmf1
mutants across six independent experiments. WT was present in all six
experiments, mmf1-1 in three and mmf1-2 in five. Values were
z-score--transformed within experiment to account for between-experiment
variation. Each point is an individual plant grown in its own pot,
i.e.~a biological replicate; number of replicates is given below each
plot. Boxplots show the median (centre line) and the 25th and 75th
percentiles (box bounds); whiskers extend to the most extreme value
within 1.5 × the interquartile range of the nearest box bound, and all
plants are plotted as points. The statistics above the plots are
estimated marginal mean (EMM) contrasts of the BX effect within each
genotype (ΔEMM = BXplus − BXminus), computed from a linear mixed-effects
model with experiment as a random intercept to account for variation
between experiments. Each ΔEMM is a two-sided t-test on the estimated
marginal means (emmeans). ΔEMM is therefore expressed in the units of
the measurement, not in z-score units. (b) Summary of the same linear
mixed-effects model in (a). For each genotype, the BX effect is reported
as the EMM contrast Δ(BXplus − BXminus) with unadjusted two-sided P
values. The table also reports interaction contrasts relative to WT (ΔΔ
vs WT = Δ\_mutant − Δ\_WT), which test whether the BX response of each
mutant differs from that of WT; these two comparisons form one family
per measure and are two-sided t-tests with Holm adjusted P values. Data
are from Dataset 4.

\subsubsection{}\label{section}

\subsubsection{Plot combined mmf1 rosette
area}\label{plot-combined-mmf1-rosette-area}

\begin{figure}

\centering{

\pandocbounded{\includegraphics[keepaspectratio]{04_Mutant_screening_files/figure-pdf/fig-mmf1-area-combined-1.pdf}}

}

\caption{\label{fig-mmf1-area-combined}Extended Data Figure 5a (rosette
area) \textbar{} Growth feedback in a BXplus microbiome is consistently
lost in mmf1 mutants. (a and b) Repeated growth feedbacks (Experiment
III - V, VIII, IX and X) of Col-0 wild type (WT) and two mmf1 T-DNA
insertion mutants (At3g44630) grown on the BXminus and BXplus soil
microbiomes. Feedback experiments were performed on
greenhouse-conditioned soil batches BS04 (Experiment III and IV), BS09
(V and X), BS10 (IX) and BS11 (VIII). (a) Z-score--transformed fresh
weight and rosette area of WT and mmf1 mutants across six independent
experiments. WT was present in all six experiments, mmf1-1 in three and
mmf1-2 in five. Values were z-score--transformed within experiment to
account for between-experiment variation. Each point is an individual
plant grown in its own pot, i.e.~a biological replicate; number of
replicates is given below each plot. Boxplots show the median (centre
line) and the 25th and 75th percentiles (box bounds); whiskers extend to
the most extreme value within 1.5 × the interquartile range of the
nearest box bound, and all plants are plotted as points. The statistics
above the plots are estimated marginal mean (EMM) contrasts of the BX
effect within each genotype (ΔEMM = BXplus − BXminus), computed from a
linear mixed-effects model with experiment as a random intercept to
account for variation between experiments. Each ΔEMM is a two-sided
t-test on the estimated marginal means (emmeans). ΔEMM is therefore
expressed in the units of the measurement, not in z-score units. (b)
Summary of the same linear mixed-effects model in (a). For each
genotype, the BX effect is reported as the EMM contrast Δ(BXplus −
BXminus) with unadjusted two-sided P values. The table also reports
interaction contrasts relative to WT (ΔΔ vs WT = Δ\_mutant − Δ\_WT),
which test whether the BX response of each mutant differs from that of
WT; these two comparisons form one family per measure and are two-sided
t-tests with Holm adjusted P values. Data are from Dataset 4.}

\end{figure}%

\#Export Source Data

\section{Export selected datasets}\label{export-selected-datasets}

\section{\texorpdfstring{\textbf{Session
information}}{Session information}}\label{session-information}

\begin{verbatim}
R version: R version 4.5.1 (2025-06-13) 
\end{verbatim}

\begin{verbatim}
Attached packages:
\end{verbatim}

\begin{verbatim}
emmeans - 2.0.3 
lmerTest - 3.2-1 
lme4 - 2.0-6 
Matrix - 1.7-5 
ggtext - 0.1.2 
cowplot - 1.2.0 
ggpubr - 1.0.0 
ggbeeswarm - 0.7.3 
writexl - 1.5.4 
openxlsx - 4.2.8.1 
knitr - 1.52 
readxl - 1.5.0 
lubridate - 1.9.5 
forcats - 1.0.1 
stringr - 1.6.0 
dplyr - 1.2.1 
purrr - 1.2.2 
readr - 2.2.0 
tidyr - 1.3.2 
tibble - 3.3.1 
ggplot2 - 4.0.3 
tidyverse - 2.0.0 
\end{verbatim}




\end{document}
